% Chapter 4 -- Rubric: PO3 (3.3, 3.4), PO5 (5.1)
\chapter{System Design and Implementation}
\label{ch:design}

\section{Functional Design and Sub-problem Partitioning}
\label{sec:architecture}

To tackle the complexities of asymmetric multi-turn dialogue, the system is designed as a modular, multi-tier computational pipeline. The overall functional architecture decomposes the problem into five primary subsystems:
\begin{enumerate}
  \item \textbf{Data Curation and Ground-Truth Engineering:} Ingests multi-turn conversational corpora (P4G and synthetic e-commerce negotiations), enforces strict semantic labeling protocols across human and automated passes, and constructs verified, leakage-free evaluation splits.
  \item \textbf{Representation and Utterance Encoding Subsystem:} Processes raw turn tokens using pre-trained contextual language models (RoBERTa, DeBERTa-v3, TOD-BERT) through batched tensor operations, producing dense contextual utterance representations while managing hardware memory footprints.
  \item \textbf{Turn-Level Rhetorical Strategy Attribution Subsystem:} Ingests Persuader turns along with causal historical context, routing representations through a 19-member multi-task model pool trained under Asymmetric Loss to predict an 11-category primary objective with 41 auxiliary strategies, unified by a calibrated decision layer.
  \item \textbf{Dialogue-Level Intention and Trajectory Tracking Subsystem (\textit{Antiphony}):} Disentangles conversational turns into role-pure Persuader and Decider streams, executes directional cross-attention, models non-monotonic stance progression via a commitment-trajectory GRU read at the true final decider turn, and outputs calibrated binary commitment predictions.
  \item \textbf{Cross-Domain Generalization and Diagnostic Ablation Subsystem (Generic Architecture Testing):} Formulates a systematic diagnostic framework spanning four input-ablation operators, trajectory destabilization via adversarial stance reversals, tokenizer truncation directionality guards, and parametric skew-stress profiling to establish whether Antiphony functions as a domain-agnostic inductive bias for asymmetric multi-turn dialogue.
\end{enumerate}

\begin{figure}[htbp]
  \centering
  \resizebox{\textwidth}{!}{% End-to-end framework: corpora -> shared representation -> two model streams -> evaluation.
\begin{tikzpicture}[x=1cm,y=1cm]

% ---------------- group backdrops -------------------------------------------------
\draw[rounded corners=6pt, draw=NeuBord!70, fill=NeuFill!55, line width=0.7pt] (0,0.1) rectangle (3.5,-10.7);
\draw[rounded corners=6pt, draw=NeuBord!70, fill=NeuFill!55, line width=0.7pt] (4.05,0.1) rectangle (7.15,-10.7);
\draw[rounded corners=6pt, draw=XatBord!70, fill=XatFill!45, line width=0.7pt] (7.7,0.1) rectangle (12.3,-5.15);
\draw[rounded corners=6pt, draw=AmbBord!70, fill=AmbFill!55, line width=0.7pt] (7.7,-5.55) rectangle (12.3,-10.7);
\draw[rounded corners=6pt, draw=TrjBord!70, fill=TrjFill!45, line width=0.7pt] (12.85,0.1) rectangle (16,-10.7);

\node[ttl] at (1.75,-0.4) {1  Corpora};
\node[ttl] at (5.6,-0.4) {2  Shared representation};
\node[ttl, text=XatBord] at (10.0,-0.4) {3a  Intention stream};
\node[ttl, text=AmbBord] at (10.0,-6.05) {3b  Strategy stream};
\node[ttl, text=TrjBord] at (14.43,-0.4) {4  Evaluation};

% ---------------- 1 corpora -------------------------------------------------------
\node[neu, text width=2.95cm, minimum height=1.85cm] (c1) at (1.75,-1.9)
  {\textbf{Persuasion for Good}\\[1pt]\scriptsize 1{,}017 dialogues\\ 10{,}600 Persuader turns\\ donation asks for Save the Children};
\node[xat, text width=2.95cm, minimum height=1.85cm] (c2) at (1.75,-4.15)
  {\textbf{Donation-intent labels}\\[1pt]\scriptsize three annotation passes\\ final split 694 yes / 323 no};
\node[amb, text width=2.95cm, minimum height=1.85cm] (c3) at (1.75,-6.4)
  {\textbf{Strategy labels}\\[1pt]\scriptsize 41 strategies, 11 categories\\ 1.68 labels per turn\\ 48.9\% multi-label turns};
\node[trj, text width=2.95cm, minimum height=1.85cm] (c4) at (1.75,-8.65)
  {\textbf{E-commerce benchmark}\\[1pt]\scriptsize 500 canonical + 200 reversal dialogues\\ three-phase synthesis};
\draw[flow] (c1) -- (c2);
\draw[flow] (c2) -- (c3);

% ---------------- 2 shared representation ------------------------------------------
\node[neu, text width=2.7cm, minimum height=2.6cm] (r1) at (5.6,-2.15)
  {\textbf{Pretrained encoders}\\[1pt]\scriptsize RoBERTa-base, DeBERTa-v3-base, TOD-BERT for intention\\ RoBERTa-base/large, ELECTRA-base for strategy};
\node[neu, text width=2.7cm, minimum height=2.0cm] (r2) at (5.6,-5.05)
  {\textbf{Role signal}\\[1pt]\scriptsize {[}Persuader{]} and {[}Persuadee{]} markers, role embeddings on every turn};
\node[neu, text width=2.7cm, minimum height=2.6cm] (r3) at (5.6,-8.15)
  {\textbf{Input policy}\\[1pt]\scriptsize turn-by-turn for intention\\ five turns of context and left truncation for flat inputs};
\draw[flow] (r1) -- (r2); \draw[flow] (r2) -- (r3);

% ---------------- 3a intention stream -----------------------------------------------
\node[xat, text width=4.0cm, minimum height=2.6cm] (m1) at (10.0,-2.35)
  {\textbf{AntiphonyClassifier}\\[1pt]\scriptsize Role-pure dual-stream cross-attention\\ Persuadee-to-Persuader attention\\ Trajectory GRU and attention pooling};
\node[fin, minimum width=2.6cm] (o1) at (10.0,-4.4) {\scriptsize commitment: yes / no};
\draw[flowX] (m1) -- (o1);

% ---------------- 3b strategy stream ------------------------------------------------
\node[amb, text width=4.0cm, minimum height=1.2cm] (s1) at (10.0,-6.95)
  {\textbf{Multi-Task Strategy Model}\\[1pt]\scriptsize 11 categories main, 41 strategies auxiliary};
\node[amb, text width=4.0cm, minimum height=1.0cm] (s2) at (10.0,-8.2)
  {\scriptsize 19-member pool: fine-tuned encoders, few-shot, TF-IDF and zero-shot members};
\node[amb, text width=1.85cm, minimum height=0.95cm, draw=AmbBord, fill=white] (s3a) at (8.95,-9.8)
  {\scriptsize \textbf{Calibrated blend}\\ Selected ($q=2$)};
\node[amb, text width=1.85cm, minimum height=0.95cm, draw=AmbBord, fill=white, dashed] (s3b) at (11.05,-9.8)
  {\scriptsize \textbf{XGBoost stacker}\\ Evaluated ($0.551$)};
\draw[flow] (s1) -- (s2);
\draw[flow] (s2.south) -- ++(0,-0.2) -| (s3a.north);
\draw[flow, dashed] (s2.south) -- ++(0,-0.2) -| (s3b.north);

% ---------------- 4 evaluation -----------------------------------------------------
\node[trj, text width=2.7cm, minimum height=1.9cm] (e1) at (14.43,-1.6)
  {\textbf{Intention benchmark}\\[1pt]\scriptsize Macro-F1 on 153 test dialogues across Pass 1--3 splits};
\node[trj, text width=2.7cm, minimum height=2.1cm] (e2) at (14.43,-3.85)
  {\textbf{Adversarial transfer}\\[1pt]\scriptsize four input ablations, sufficiency gap $\Delta_{\mathrm{suff}}$, skew stress};
\node[trj, text width=2.7cm, minimum height=1.9cm] (e3) at (14.43,-6.2)
  {\textbf{Strategy benchmark}\\[1pt]\scriptsize micro and macro-F1 against the 0.79 annotation ceiling};
\node[trj, text width=2.7cm, minimum height=2.1cm] (e4) at (14.43,-8.6)
  {\textbf{Forensic checks}\\[1pt]\scriptsize truncation side, marker corruption, merchant-side leakage};

% ---------------- inter-group arrows -----------------------------------------------
\draw[flow, line width=1.3pt] (3.5,-2.15) -- (4.05,-2.15);
\draw[flow, line width=1.3pt] (3.5,-8.15) -- (4.05,-8.15);
\draw[flow, line width=1.3pt] (7.15,-2.6) -- (7.7,-2.6);
\draw[flow, line width=1.3pt] (7.15,-8.0) -- (7.7,-8.0);
\draw[flow, line width=1.3pt] (12.3,-2.6) -- (12.85,-2.6);
\draw[flow, line width=1.3pt] (12.3,-8.0) -- (12.85,-8.0);

\end{tikzpicture}
}
  \caption{Overall functional system architecture, illustrating data ingestion across three corpora, shared utterance encoding, the dual-stream Antiphony intention network, the 19-member multi-task strategy classification framework with calibrated decision blending, and the diagnostic ablation evaluation framework.}
  \label{fig:system_architecture}
\end{figure}

\Cref{tab:subsystem-responsibilities} enumerates the functional modules, their inputs, their outputs, and their architectural responsibilities.

\begin{table}[htbp]
  \centering
  \small
  \caption{Modular subsystem decomposition and responsibilities}
  \label{tab:subsystem-responsibilities}
  \begin{tabularx}{\textwidth}{l>{\hsize=0.8\hsize}Y>{\hsize=0.8\hsize}Y>{\hsize=1.4\hsize}Y}
    \toprule
    \textbf{Subsystem} & \textbf{Input} & \textbf{Output} & \textbf{Core Engineering Responsibility} \\
    \midrule
    \textbf{Annotation \& Ground Truth} & Raw text transcripts ($N=1{,}017$, $10{,}600$ turns) & Grounded labels ($y \in \{\text{yes},\text{no}\}$, $\mathbf{y}_t \subseteq \mathcal{S}_{41}$) & Resolves metadata divergence via 3-pass human annotation; enforces 4-tier decision protocol. \\
    \midrule
    \rowcolor{lightbg}
    \textbf{Utterance Encoder} & Batched token IDs $[B \times U, L]$ & Turn matrices $\mathbf{U} \in \mathbb{R}^{B \times U \times d}$ & Executes single batched forward pass; performs masked mean-pooling; enforces memory guards. \\
    \midrule
    \textbf{Strategy Classifier} & Persuader turn $u_t + \mathcal{H}_t$ (5-turn context) & Predicted categories $\hat{\mathbf{c}}_t$ \& strategies $\hat{\mathbf{y}}_t$ & Multi-task learning under Asymmetric Loss; ensembles 19 members via calibrated power-mean blend. \\
    \midrule
    \rowcolor{lightbg}
    \textbf{Antiphony Intention Network} & Full dialogue turn vectors $\mathbf{U}$ & Commitment probability $P(y = \text{yes} \mid \mathcal{D})$ & Role disentanglement ($R=20$); directional cross-attention; trajectory GRU tracking; dual pooling fusion. \\
    \midrule
    \textbf{Diagnostic Ablation Harness} & Full and ablated dialogue streams $\mathbf{H}_M, \mathbf{H}_B$ & Counterfactual ablations (LTR, LTR+M, FHO, PRO), $\Delta_{\text{suff}}$, stress curves & Enforces left-truncation invariant; isolates early merchant accommodation signals; verifies domain-agnostic trajectory necessity. \\
    \bottomrule
  \end{tabularx}
\end{table}

\section{Persuasion Strategy Annotation and Taxonomy}
\label{sec:strategy-annotation}

Modeling the rhetorical tactics that influence conversational outcomes requires an empirically validated, fine-grained taxonomy and a standardized annotation protocol. We establish a multi-label persuasion strategy benchmark encompassing the first $10{,}600$ Persuader turns from $1{,}017$ dialogues in the P4G corpus \cite{wang2019persuasion}.

\subsection{Hierarchical Taxonomy Formulation}
Adapted from social influence theory \cite{cialdini2001influence}, compliance-gaining literature \cite{freedman1966compliance,cialdini1975reciprocal}, and pragmatic dialogue analysis, our taxonomy establishes a two-tier hierarchy: $41$ fine-grained rhetorical strategies nested deterministically under $11$ functional communicative categories (structured in a 4, 4, 4, 4, 3, 4, 4, 3, 3, 4, 4 layout):
\begin{enumerate}[itemsep=2pt,parsep=0pt]
  \item \textbf{Rational Appeal (4 strategies):} \sloppy \texttt{logical\_appeal}, \texttt{evidence\_and\_statistics}, \texttt{cost\_benefit\_framing}, \texttt{feasibility\_and\_ease}.
  \item \textbf{Emotional Appeal (4 strategies):} \sloppy \texttt{emotion\_appeal}, \texttt{guilt\_induction}, \texttt{empathy\_and\_perspective\_taking}, \texttt{hope\_and\_positive\_impact}.
  \item \textbf{Credibility Appeal (4 strategies):} \sloppy \texttt{organizational\_credibility}, \texttt{transparency\_and\_accountability}, \texttt{source\_citation}, \texttt{personal\_credibility}.
  \item \textbf{Social Influence (4 strategies):} \sloppy \texttt{social\_proof}, \texttt{self\_modeling}, \texttt{in\_group\_appeal}, \texttt{authority\_endorsement}.
  \item \textbf{Reciprocity and Exchange (3 strategies):} \sloppy \texttt{reciprocity}, \texttt{donor\_benefit}, \texttt{gratitude\_and\_appreciation}.
  \item \textbf{Commitment and Consistency (4 strategies):} \sloppy \texttt{foot\_in\_the\_door}, \texttt{door\_in\_the\_face}, \texttt{value\_consistency\_appeal}, \texttt{incremental\_ask}.
  \item \textbf{Framing and Presentation (4 strategies):} \sloppy \texttt{anchoring}, \texttt{minimization\_framing}, \texttt{loss\_versus\_gain\_framing}, \texttt{comparison\_framing}.
  \item \textbf{Urgency and Scarcity (3 strategies):} \sloppy \texttt{urgency\_appeal}, \texttt{scarcity\_appeal}, \texttt{deadline\_pressure}.
  \item \textbf{Threat and Pressure (3 strategies):} \sloppy \texttt{negative\_consequence\_warning}, \texttt{persistent\_repetition}, \texttt{obligation\_pressure}.
  \item \textbf{Relational and Interactive (4 strategies):} \sloppy \texttt{rapport\_building}, \texttt{personal\_story}, \texttt{personal\_related\_inquiry}, \texttt{source\_related\_inquiry}.
  \item \textbf{Information Provision (4 strategies):} \sloppy \texttt{organization\_information}, \texttt{donation\_procedure\_information}, \texttt{task\_clarification}, \texttt{impact\_information}.
\end{enumerate}

\subsection{Contextual Annotation Pipeline and Decision Protocol}
Annotating isolated sentences distorts rhetorical meaning; an utterance such as ``Would you be willing to spare 50 cents?'' functions as an \texttt{incremental\_ask} only if preceded by a larger rejected request. Therefore, our pipeline ingests each Persuader utterance within a 5-turn preceding conversational context window ($[u_{t-5}, \dots, u_{t-1}]$) maintaining true speaker alternation.

\begin{figure}[htbp]
  \centering
  \resizebox{\textwidth}{!}{% Multi-label strategy annotation pipeline: 10,600 Persuader turns, 41 strategies, 11 categories.
% Layout: data preparation runs down the left column, then the labelling run flows
% left to right along the bottom row. The supporting groups sit directly above the
% step they act on: decision tests -> prompt, human checkpoints -> labelling and
% validation, merged corpus -> resulting corpus.
% Boxes are stacked with relative positioning and every group backdrop is computed
% with `fit`, so boxes cannot overlap or extend past their group.
% Width is about 16.4 cm; wrap in \resizebox{\textwidth}{!}{...} if the text block is narrower.
\begin{tikzpicture}[x=1cm,y=1cm, node distance=0.32cm,
  gttl/.style={font=\footnotesize\bfseries, anchor=north, inner sep=2pt},
  band/.style={rounded corners=6pt, line width=0.7pt, inner sep=0.17cm},
  card/.style={blk, fill=white, text width=2.4cm},
  prep/.style={neu, text width=2.4cm},
  stp/.style={text width=2.45cm, minimum height=1.9cm},
  mainflow/.style={flow, line width=1.1pt},
]

% column centres: 5 columns, 2.9 cm wide, 0.475 cm gutters
\def\cA{1.45}\def\cB{4.825}\def\cC{8.2}\def\cD{11.575}\def\cE{14.95}
\coordinate (xA) at (\cA,0); \coordinate (xC) at (\cC,0);
\coordinate (xD) at (\cD,0); \coordinate (xE) at (\cE,0);

% ======================= column 1: data preparation (top to bottom) ================
\node[gttl] (gA) at (\cA,0) {Data preparation};
\node[prep, below=0.12cm of gA] (a1)
  {\textbf{Persuasion for Good}\\[1pt]\scriptsize 1{,}017 two-party dialogues};
\node[prep, below=0.42cm of a1] (a2)
  {\textbf{Persuader turn table}\\[1pt]\scriptsize 10{,}600 turns in canonical chronological order};
\node[prep, below=0.42cm of a2] (a3)
  {\textbf{Contiguous batches}\\[1pt]\scriptsize 200 turns each, unshuffled, dialogues kept together};

% ======================= column 2: four decision tests ==============================
\node[gttl, text=XatBord] (gB) at (\cB,0) {Four decision tests};
\node[xat, card, below=0.12cm of gB] (t1)
  {\textbf{1\enspace Removal test}\\[1pt]\scriptsize a label needs evidence that survives deleting the others};
\node[xat, card, below=of t1] (t2)
  {\textbf{2\enspace Specific over general}\\[1pt]\scriptsize the fine strategy replaces its parent concept};
\node[xat, card, below=of t2] (t3)
  {\textbf{3\enspace Function over keyword}\\[1pt]\scriptsize pragmatic role in context, not trigger words};
\node[xat, card, below=of t3] (t4)
  {\textbf{4\enspace Density cap}\\[1pt]\scriptsize one to three labels typical, empty set allowed, six at most};

% ======================= columns 3-4: human evaluation checkpoints ==================
\coordinate (mCD) at ($(\cC,0)!0.5!(\cD,0)$);
\node[gttl, text=AmbBord] (gC) at (mCD) {Human evaluation checkpoints};
\node[amb, card, anchor=north] (h1) at ($(\cC,0)+(0,-0.62)$)
  {\textbf{Smoke test}\\[1pt]\scriptsize the first batch of 200 is labelled alone; the team checks every label and signs off before the run continues};
\node[amb, card, anchor=north] (h2) at ($(\cD,0)+(0,-0.62)$)
  {\textbf{Spot checks}\\[1pt]\scriptsize sampled slices of later batches are audited for drift, especially rare strategies and turns with four or more labels};
\node[nt, text width=2.55cm, anchor=north] (note) at ($(mCD |- h2.south)+(0,-0.3)$)
  {labelled by a frontier LLM agent under these checkpoints; no turn was labelled twice};

% ======================= column 5: resulting corpus =================================
\node[gttl] (gE) at (\cE,0) {Resulting corpus};
\node[neu, card, anchor=north] (st) at ($(\cE,0)+(0,-0.62)$)
  {\scriptsize mean 1.68 labels per turn\\[2pt] 48.9\% multi-label turns\\[2pt] 6.7\% empty sets\\[2pt] all 41 strategies and all 11 categories occur};

% ======================= top group backdrops (fitted) ===============================
% Groups share a top edge and end at their own natural height.
\begin{scope}[on background layer]
  \node[band, draw=XatBord!70, fill=XatFill!40, fit=(gB)(t1)(t4)] (GB) {};
  \node[band, draw=AmbBord!70, fill=AmbFill!45,
        fit=(gC)(h1)(h2)(note)(GB.north -| h1.west)] (GC) {};
  \node[band, draw=NeuBord!70, fill=NeuFill!60, fit=(gE)(st)] (GE) {};
\end{scope}

% ======================= main row: the labelling run ================================
\node[xat, stp, anchor=north] (b1) at ($(GB.south)+(0,-0.75)$)
  {\textbf{Prompt assembly}\\[1pt]\scriptsize taxonomy of 11 + 41 labels, four decision tests, five gold few-shot examples};
\node[amb, stp] (b2) at (xC |- b1)
  {\textbf{LLM agent labels a batch}\\[1pt]\scriptsize JSON per turn: strategies, categories, n\_labels, flags, rationale};
\node[trj, stp] (b3) at (xD |- b1)
  {\textbf{Algorithmic validation}\\[1pt]\scriptsize closed-world taxonomy, category invariant, at most six labels};
\node[fin, stp] (b4) at (xE |- b1)
  {\textbf{Merged corpus}\\[1pt]\scriptsize JSONL, CSV, and a $10{,}600\times56$ binary matrix};
\node[prep, minimum height=1.9cm] (a4) at (xA |- b1)
  {\textbf{Five-turn context window}\\[1pt]\scriptsize five preceding turns, alternating roles};

\begin{scope}[on background layer]
  \node[band, draw=NeuBord!70, fill=NeuFill!55, fit=(gA)(a1)(a2)(a3)(a4)] (GA) {};
\end{scope}

% ======================= arrows =====================================================
% data preparation, top to bottom
\draw[flow] (a1) -- (a2); \draw[flow] (a2) -- (a3); \draw[flow] (a3) -- (a4);
% main labelling run, left to right
\draw[mainflow] (a4) -- (b1);
\draw[mainflow] (b1) -- (b2);
\draw[mainflow] (b2) -- (b3);
\draw[mainflow] (b3) -- (b4);
% decision tests are written into the prompt
\draw[flowX] (t1) -- (t2); \draw[flowX] (t2) -- (t3); \draw[flowX] (t3) -- (t4);
\draw[flowX, dsh] (t4.south) -- (b1.north);
% human checkpoints act on labelling and validation
\draw[flow, dsh, draw=AmbBord] (h1.south) -- (b2.north);
\draw[flow, dsh, draw=AmbBord] (h2.south) -- (b3.north);
% merged corpus summarised
\draw[flow, dsh] (b4.north) -- (st.south);

\end{tikzpicture}
}
  \caption{Persuasion strategy annotation pipeline, depicting contextual batching (200 turns per batch, 5-turn history), the 4-tier decision protocol, LLM agent annotation under human checkpoints, automated schema validation, and multi-label compilation into \texttt{multilabel\_merged.jsonl}.}
  \label{fig:strategy_annotation_pipeline}
\end{figure}

To eliminate subjective drift and enforce consistency across segments, annotators followed a mandatory \textbf{Four-Tier Decision Protocol} applied in fixed sequential order:
\begin{enumerate}
  \item \textbf{The Removal Test:} A strategy label $\mathbf{s}$ is valid if and only if removing the corresponding clause or rhetorical move noticeably alters the persuasive force or communicative intent of the turn.
  \item \textbf{Specific-over-General Precedence:} If an utterance satisfies both a general category definition (e.g., \texttt{logical\_appeal}) and a highly specific operational tactic (e.g., \texttt{evidence\_and\_statistics} or \texttt{cost\_benefit\_framing}), the specific tactic must be coded, avoiding redundant double-labeling.
  \item \textbf{Pragmatic Function over Keyword Matching:} Labels are assigned based on contextual communicative function rather than literal lexical cues. A rhetorical question (``Don't you want to help a dying child?'') is coded as \texttt{guilt\_induction} and \texttt{emotion\_appeal}, never as an informational inquiry.
  \item \textbf{Realistic Label Density Cap:} While multi-label combinations are expected in complex turns, humans rarely deploy more than three distinct rhetorical tactics simultaneously. The pipeline enforces a typical density of 1 to 3 labels, permits an explicit empty set ($\mathbf{y}_t = \emptyset$) for purely conversational or logistical acknowledgments, and enforces a hard ceiling of $|\mathbf{y}_t| \le 6$.
\end{enumerate}

Annotation was conducted in 200-turn unshuffled batches using an automated frontier LLM agent (Claude 3.5 Sonnet) driven with few-shot calibration examples, operating under human checkpoints. An automated validation harness (\texttt{pipeline.py}) verified taxonomy compliance, recomputed parent categories from strategies, enforced the label count ceiling, and compiled the verified dataset into \texttt{multilabel\_merged.jsonl} ($10{,}600 \times 56$ wide format with 41 strategy and 11 category flags).

\begin{terminalbox}[Execution Trace: Strategy Annotation Schema Validation (\texttt{pipeline.py})]
[+] Ingesting annotated chunks from \texttt{data/raw\_batches/} (10,600 Persuader turns) \\ \relax
[+] Running 4-tier decision constraint harness: \\ \relax
\hspace*{1.5em}- Verifying max strategy density cap: $\max(|\mathbf{y}_t|) \le 6$ \dots\ \textbf{PASS} (mean = 1.68) \\ \relax
\hspace*{1.5em}- Validating deterministic category mapping $\mathcal{M}: \mathcal{S}_{41} \to \mathcal{C}_{11}$ \dots\ \textbf{PASS} (100\% consistent) \\ \relax
\hspace*{1.5em}- Enforcing closed-world vocabulary constraints \dots\ \textbf{PASS} (0 unknown tokens) \\ \relax
\hspace*{1.5em}- Checking speaker alternation invariants \dots\ \textbf{PASS} ($\tau_{\mathcal{P}} = 2t - 1$) \\ \relax
[*] Merging verified batches into \texttt{multilabel\_merged.jsonl} ($10{,}600 \times 56$ matrix) \\ \relax
[+] DATASET AUDIT COMPLETE: 10,600 records serialized, 0 schema violations.
\end{terminalbox}

\subsection{Empirical Corpus Statistics and Rater Effects}
The resulting strategy corpus encompasses $10{,}600$ Persuader turns across $1{,}017$ dialogues (mean $10.42$ Persuader turns per dialogue, range 4 to 15). The dataset exhibits an average of $1.679$ labels per turn ($\text{sd} = 1.061$, median $1.0$). Exactly $707$ turns ($6.7\%$) carry no strategy label (the empty set $\emptyset$), $4{,}707$ turns ($44.4\%$) carry exactly 1 label, and $5{,}182$ turns ($48.9\%$) are multi-label ($\ge 2$ strategies), with 78 turns reaching the maximum cap of 6 labels.

\Cref{tab:category-prevalence} reports the empirical prevalence of all 11 communicative categories, recomputed directly from the data-verified \texttt{multilabel\_merged.jsonl}.

\begin{table}[htbp]
  \centering
  \small
  \caption{Empirical prevalence of communicative categories across 10,600 Persuader turns}
  \label{tab:category-prevalence}
  \begin{tabularx}{\textwidth}{l>{\raggedleft\arraybackslash}X>{\raggedleft\arraybackslash}X}
    \toprule
    \textbf{Communicative Category} & \textbf{Active Turns} & \textbf{Prevalence (\%)} \\
    \midrule
    Relational and Interactive & 4,897 & 46.2\% \\
    \rowcolor{lightbg}
    Information Provision & 3,823 & 36.1\% \\
    Emotional Appeal & 1,645 & 15.5\% \\
    \rowcolor{lightbg}
    Social Influence & 1,263 & 11.9\% \\
    Reciprocity and Exchange & 1,078 & 10.2\% \\
    \rowcolor{lightbg}
    Framing and Presentation & 1,004 & 9.5\% \\
    Credibility Appeal & 896 & 8.5\% \\
    \rowcolor{lightbg}
    Rational Appeal & 751 & 7.1\% \\
    Threat and Pressure & 249 & 2.3\% \\
    \rowcolor{lightbg}
    Commitment and Consistency & 169 & 1.6\% \\
    Urgency and Scarcity & 163 & 1.5\% \\
    \bottomrule
  \end{tabularx}
\end{table}

\textbf{Annotator Granularity Effects:} The corpus was annotated across four distinct operational segments under human checkpoints: Segment 1 (turns 1--2,000; Track A), Segment 2 (turns 2,001--4,000 and 6,001--7,000; Track B), Segment 3 (turns 4,001--6,000; Calibration Track), and Segment 4 (turns 7,001--10,600; Track C). Analysis reveals significant variance in annotation granularity: Track A averaged $2.384$ labels per turn ($9.0\%$ empty sets, $66.8\%$ multi-label), whereas Track C averaged $1.240$ labels per turn ($13.2\%$ empty sets, $30.8\%$ multi-label). This variance represents an operational rater effect in label granularity rather than an intrinsic linguistic shift in the underlying conversations.

\subsection{Derivation of the Cross-Annotator Agreement Ceiling}
Because individual turns were annotated once without global multi-coder duplication, inter-annotator agreement cannot be computed via standard Cohen's kappa. Instead, we derive the empirical performance ceiling from near-duplicate turns across the corpus. On identical or near-identical text pairs (cosine similarity $\ge 0.95$):
\begin{itemize}
  \item When the \textit{same} annotator labels near-identical turns, their self-agreement reaches an F1 of $0.871$.
  \item When \textit{different} annotators label near-identical turns, their cross-agreement drops to an F1 of $0.657$.
\end{itemize}
Because a production model never receives annotator identity at inference time, it must generalize across annotators. Applying consensus prediction formulations under this observed rater noise establishes a practical upper performance ceiling of \textbf{$\approx 0.79$ micro-F1}. All strategy model evaluations must be interpreted relative to this $0.79$ ceiling, where a score of $0.70$ captures nearly $89\%$ of theoretically achievable human consensus.

\section{Conversational Intention Ground Truth Generation}
\label{sec:intention-annotation}

\subsection{The Ground-Truth Divergence Problem}
The original P4G benchmark \cite{wang2019persuasion} recorded numeric donation values reported by Mechanical Turk workers. However, systematic manual auditing revealed severe discrepancies between tabular records and actual conversational text:
\begin{itemize}
  \item In Dialogue \#124, the Persuadee explicitly confirms: `\texttt{We have you down for a \$10 donation. Is that correct?'' / }`Yes that is'', yet the recorded donation amount is \$0.
  \item In Dialogue \#842, the Persuadee states: ``I will absolutely look to donating something later'', while both numeric donation fields show \$0.
\end{itemize}
Training machine learning classifiers on unverified tabular values causes models to learn spurious logging noise rather than genuine conversational commitments.

\begin{figure}[htbp]
  \centering
  \resizebox{\textwidth}{!}{% Three-pass engineering of the donation-intent labels (1,017 dialogues).
% Bar lengths are proportional: 4.6 cm = 1,017 dialogues.
\begin{tikzpicture}[x=1cm,y=1cm]

% ---------------- frames and headers ------------------------------------------------
\foreach \x/\c/\t in {0/94A3B8/{First pass: manual dual-field labels}, 5.5/64748B/{Second pass: plain binary}, 11/334155/{Third pass: strict commitment}}{
  \definecolor{stg}{HTML}{\c}
  \draw[rounded corners=5pt, draw=stg, fill=white, line width=0.9pt] (\x,0.05) rectangle ++(5,-9.65);
  \fill[stg, rounded corners=5pt] (\x,0.05) rectangle ++(5,-0.75);
  \fill[stg] (\x,-0.4) rectangle ++(5,-0.3);
  \node[text=white, font=\footnotesize\bfseries] at (\x+2.5,-0.33) {\t};
}
\foreach \x in {5,10.5}{ \draw[flow, line width=1.3pt] (\x-0.02,-0.35) -- ++(0.54,0); }

% ---------------- body text ---------------------------------------------------------------
\node[anchor=north west, text width=4.7cm, font=\scriptsize, align=left] at (0.15,-0.95)
  {$\bullet$ pilot of 200 dialogues split five ways, 50 each, with an annotator column for a drift check\\
   $\bullet$ then about 200 dialogues per member, pooled to 1{,}017\\
   $\bullet$ fields: \texttt{binary\_label} and \texttt{modifier} (none, deferred, conditional)\\
   $\bullet$ motivation: the corpus donation fields disagree with the transcript};
\node[anchor=north west, text width=4.7cm, font=\scriptsize, align=left] at (5.65,-0.95)
  {$\bullet$ supervisor directive: drop the modifier and focus on the binary decision\\
   $\bullet$ the rarest modifier class had 9 training and 3 test examples\\
   $\bullet$ \texttt{binary\_label} is left exactly as annotated\\
   $\bullet$ basis of Decoupled Baseline and Antiphony initial evaluations};
\node[anchor=north west, text width=4.7cm, font=\scriptsize, align=left] at (11.15,-0.95)
  {$\bullet$ any hesitation, condition, or deferral now counts as \emph{no}\\
   $\bullet$ ``I'll donate \$1 now, and I might give more later'' becomes \emph{no}\\
   $\bullet$ built in two steps, shown below\\
   $\bullet$ basis of terminal benchmarking for Decoupled and Antiphony models};

% ---------------- bars: first pass ------------------------------------------------------------
\begin{scope}[shift={(0.2,0)}]
  \node[font=\tiny, anchor=west, text=SlateSoft] at (0,-3.65) {commitment label};
  \fill[TrjMain] (0,-3.8) rectangle (3.338,-4.2); \fill[AmbMain] (3.338,-3.8) rectangle (4.6,-4.2);
  \node[font=\scriptsize\bfseries, text=white] at (1.67,-4.0) {yes 738 (72.6\%)};
  \node[font=\scriptsize\bfseries, text=white] at (3.97,-4.0) {no 279};
  \node[font=\tiny, anchor=west, text=SlateSoft] at (0,-4.5) {modifier among the yes dialogues};
  \fill[TrjMain!55] (0,-4.65) rectangle (2.972,-5.0);
  \fill[XatMain] (2.972,-4.65) rectangle (3.28,-5.0);
  \fill[PerMain] (3.28,-4.65) rectangle (3.339,-5.0);
  \node[font=\scriptsize, text=white] at (1.49,-4.83) {none 657};
  \node[font=\tiny, anchor=north east, text=XatBord] at (3.45,-5.05) {deferred 68};
  \node[font=\tiny, anchor=north west, text=PerBord] at (3.52,-5.05) {conditional 13};
  \draw[XatMain, line width=0.5pt] (3.126,-5.0) -- (3.126,-5.06); \draw[PerMain, line width=0.5pt] (3.31,-5.0) -- (3.45,-5.06);
\end{scope}

% ---------------- bars: second pass -----------------------------------------------------------
\begin{scope}[shift={(5.7,0)}]
  \node[font=\tiny, anchor=west, text=SlateSoft] at (0,-3.65) {commitment label};
  \fill[TrjMain] (0,-3.8) rectangle (3.338,-4.2); \fill[AmbMain] (3.338,-3.8) rectangle (4.6,-4.2);
  \node[font=\scriptsize\bfseries, text=white] at (1.67,-4.0) {yes 738 (72.6\%)};
  \node[font=\scriptsize\bfseries, text=white] at (3.97,-4.0) {no 279};
  \node[font=\tiny, anchor=west, text=SlateSoft] at (0,-4.5) {modifier target (deprecated)};
  \draw[SlateSoft!60, line width=0.6pt, dashed, fill=black!3] (0,-4.65) rectangle (3.338,-5.0);
  \draw[red!80!black, line width=0.7pt] (0,-4.65) -- (3.338,-5.0);
  \draw[red!80!black, line width=0.7pt] (0,-5.0) -- (3.338,-4.65);
  \node[font=\tiny\bfseries, text=red!80!black, fill=white, inner sep=1.2pt] at (1.67,-4.825) {DEPRECATED HEAD};
\end{scope}

% ---------------- bars: third pass ------------------------------------------------------------
\begin{scope}[shift={(11.2,0)}]
  \node[font=\tiny, anchor=west, text=SlateSoft] at (0,-3.65) {after step 1 (deterministic remap)};
  \fill[TrjMain] (0,-3.8) rectangle (2.972,-4.2); \fill[AmbMain] (2.972,-3.8) rectangle (4.6,-4.2);
  \node[font=\scriptsize\bfseries, text=white] at (1.49,-4.0) {yes 657};
  \node[font=\scriptsize\bfseries, text=white] at (3.79,-4.0) {no 360};
  \node[font=\tiny, anchor=west, text=SlateSoft] at (0,-4.5) {final (after step 2, manual re-audit)};
  \fill[TrjMain] (0,-4.65) rectangle (3.139,-5.05); \fill[AmbMain] (3.139,-4.65) rectangle (4.6,-5.05);
  \node[font=\scriptsize\bfseries, text=white] at (1.57,-4.85) {yes 694 (68.2\%)};
  \node[font=\scriptsize\bfseries, text=white] at (3.87,-4.85) {no 323};
\end{scope}

% ---------------- third pass, the two steps --------------------------------------------------
\node[amb, fill=white, text width=4.5cm, minimum height=1.25cm, anchor=north west] (s1) at (11.15,-5.45)
  {\textbf{Step 1}\\[-1pt]\scriptsize \emph{yes} with modifier \emph{none} stays \emph{yes}; everything else becomes \emph{no}. 81 dialogues move from yes to no.};
\node[amb, fill=white, text width=4.5cm, minimum height=1.6cm, anchor=north west] (s2) at (11.15,-7.05)
  {\textbf{Step 2}\\[-1pt]\scriptsize all 81 are read in full and judged by the Persuadee's last position: 44 are confirmed \emph{no}, 37 had resolved their hedge and revert to \emph{yes}.};
\draw[flow] (s1.south) -- (s2.north);

% ---------------- what the earlier passes leave open ---------------------------------------
\node[anchor=north west, text width=4.7cm, font=\scriptsize, align=left] at (0.15,-5.75)
  {\textbf{Reliability check.} An independent second pass over 154 dialogues of the lead partition, made without sight of the first labels, agreed on 142 of them (92.21\%). The 12 disagreements clustered around four boundary cases and were written back into the guideline.};
\node[anchor=north west, text width=4.7cm, font=\scriptsize, align=left] at (5.65,-5.75)
  {\textbf{Why this version exists.} With modifier gradients removed, the shared representation no longer trades binary accuracy against a data-starved head; every binary result in the Second-Pass Benchmark comes from these labels.};

% ---------------- where each pass is used -------------------------------------------------------
\foreach \x/\t in {0.2/{First-Pass Benchmark (Dual-Head)}, 5.7/{Second-Pass Benchmark (Decoupled)}, 11.2/{Third-Pass Strict Benchmark}}{
  \node[trj, fill=TrjFill, text width=4.5cm, minimum height=0.6cm, anchor=north west] at (\x,-8.85) {\scriptsize\bfseries \t};
}

\end{tikzpicture}
}
  \caption{Three-pass intention annotation lineage, showing the initial 738/279 binary split with modifier, the second pass deprecating the modifier, and the third pass enforcing strict unconditional commitment (694 yes / 323 no via two-step re-checking).}
  \label{fig:intention_annotation_passes}
\end{figure}

\subsection{The Three-Pass Annotation Lineage}
To establish an unassailable ground truth, we executed three successive annotation passes:
\begin{enumerate}
  \item \textbf{First Pass (Dual-Field Manual Annotation):} A pilot of 200 dialogues was split five ways (50 per annotator) with an annotator-tracking column to verify calibration and prevent rater drift. Following calibration, the five team members independently annotated $\approx 200$ dialogues each, reading full transcripts and assigning \texttt{binary\_label} (\texttt{yes}/\texttt{no}) and a \texttt{modifier} (\texttt{none}, \texttt{deferred}, \texttt{conditional}). The pooled 1,017-dialogue distribution yielded: $738$ \texttt{yes} ($72.6\%$) and $279$ \texttt{no} ($27.4\%$). Among affirmative commitments, $657$ were \texttt{none}, $68$ were \texttt{deferred}, and only $13$ were \texttt{conditional} (9 in training, 3 in test).
  \item \textbf{Second Pass (Modifier Deprecation):} Modeling experiments (Chapter 5) proved that the modifier head severely degraded binary performance and suffered from severe data sparsity on the 13-example \texttt{conditional} class. Following a directive from the project supervisor, \SupervisorName, the modifier target was dropped while leaving binary labels unchanged ($738$ \texttt{yes} / $279$ \texttt{no}).
  \item \textbf{Third Pass (Strict Unconditional Commitment):} To create an unambiguous, operationally robust target, the definition of \texttt{yes} was tightened: any hesitation, conditional clause (``if'', ``maybe''), or temporal deferral (``later'') was strictly reclassified as \texttt{no}. This was executed in two steps:
    \begin{itemize}
      \item \textit{Step 1 (Deterministic Remap):} Any dialogue originally marked \texttt{yes} with modifier \texttt{none} remained \texttt{yes}; all dialogues with \texttt{deferred} or \texttt{conditional} modifiers were remapped to \texttt{no}, moving 81 dialogues from \texttt{yes} to \texttt{no} ($657$ \texttt{yes} / $360$ \texttt{no}).
      \item \textit{Step 2 (Terminal Stance Verification):} Because a modifier tag only indicated that a hedge occurred somewhere in the chat, all 81 remapped dialogues were manually re-read in full to evaluate their final state. If an early hesitation was cleanly resolved into an unconditional affirmative by the end, it was reverted to \texttt{yes}. Exactly 44 dialogues were confirmed as \texttt{no}, and 37 were reverted to \texttt{yes}.
    \end{itemize}
    The finalized third-pass ground truth contains exactly \textbf{694 \texttt{yes} ($68.2\%$) / 323 \texttt{no} ($31.8\%$)}.
\end{enumerate}

\subsection{Inter-Annotator Agreement and Boundary Cases}
An independent second-pass annotation over 154 dialogues established an inter-annotator agreement rate of \textbf{$92.21\%$} (142 of 154 identical assignments). Qualitative review of the 12 boundary disagreements revealed four recurring linguistic phenomena: (i) ambiguous temporal deferrals (``I'll donate through their website tonight''), (ii) polite deflection masquerading as consent (``Sure, that sounds like a great cause''), (iii) conditional matches (``I'll give \$1 if you give \$1''), and (iv) unconfirmed donation queries. Adjudicating these boundary cases under the strict third-pass standard eliminated residual ambiguity.

\section{Synthetic E-Commerce Negotiation Benchmark}
\label{sec:synthetic-benchmark}

\subsection{Motivation and Privacy Bottlenecks}
Evaluating cross-domain generalization from charitable persuasion to commercial negotiation is obstructed by data access barriers: private customer-merchant chat logs on platforms like WhatsApp or Facebook Messenger are legally protected under privacy statutes and commercially sensitive. To enable rigorous, leak-free evaluation, we designed a three-phase synthetic benchmark generation pipeline producing 700 fully audited negotiation dialogues.

\begin{figure}[htbp]
  \centering
  \resizebox{\textwidth}{!}{% Three-phase synthetic e-commerce negotiation benchmark: synthesis, humanisation, audit.
% Each phase is a column read top to bottom. The artefact each phase hands on sits in a
% common bottom row, so the hand-offs read left to right (drafts -> humanised -> frozen).
% Regimes A and B are parallel prompt regimes feeding one generation engine.
% Phase frames are computed with `fit`, so no box can extend past its frame.
% Width is about 16.4 cm; wrap in \resizebox{\textwidth}{!}{...} if the text block is narrower.
\begin{tikzpicture}[x=1cm,y=1cm, node distance=0.32cm,
  phead/.style={rectangle, rounded corners=4pt, fill=#1, text=white, font=\footnotesize\bfseries,
                minimum height=0.62cm, anchor=north},
  pframe/.style={rounded corners=6pt, draw=#1, line width=0.9pt, fill=white, inner sep=0.16cm},
  wide/.style={text width=4.3cm},
  outbox/.style={blk, text width=4.3cm, minimum height=1.75cm, line width=1.0pt},
  handoff/.style={flow, line width=1.3pt, draw=Slate},
]
\definecolor{PhA}{HTML}{94A3B8}\definecolor{PhB}{HTML}{64748B}\definecolor{PhC}{HTML}{334155}

% column spans: phase 1 0-5.6, phase 2 6.1-11.0, phase 3 11.5-16.4
\coordinate (c1) at (2.8,0); \coordinate (c2) at (8.55,0); \coordinate (c3) at (13.95,0);

% ======================= headers ====================================================
\node[phead=PhA, minimum width=5.6cm] (H1) at (c1) {Phase 1\enspace Persona-driven synthesis};
\node[phead=PhB, minimum width=4.9cm] (H2) at (c2) {Phase 2\enspace Organic paraphrasing};
\node[phead=PhC, minimum width=4.9cm] (H3) at (c3) {Phase 3\enspace Labelling and audit};

% ======================= phase 1 ====================================================
% two parallel prompt regimes
\node[per, text width=2.3cm, minimum height=4.0cm, anchor=north east] (rA) at ($(H1.south)+(-0.06,-0.35)$)
  {\textbf{Regime A}\\\textbf{canonical, $N{=}500$}\\[2pt]\scriptsize 8 buyer archetypes: Skeptic, LowBaller, Urgent, Ghoster, Comparator, DeliveryObsessed, TrustSeeker, BulkBuyer\\[1pt] 3 merchant personas\\[1pt] 4 product categories};
\node[dec, text width=2.3cm, minimum height=4.0cm, anchor=north west] (rB) at ($(H1.south)+(0.06,-0.35)$)
  {\textbf{Regime B}\\\textbf{reversals, $N{=}200$}\\[2pt]\scriptsize 65 resistance $\to$ agreement\\[1pt] 65 commitment $\to$ rejection\\[1pt] 70 hard monotonic controls\\[1pt] 6 buyer sentiment states and a \texttt{reversal\_turn\_id}};
% both regime boxes share a fixed height; take the lower edge to be safe
\path let \p1=(rA.south), \p2=(rB.south) in coordinate (regBot) at (0,{min(\y1,\y2)});
\node[xat, wide, text width=5.0cm, anchor=north] (eng) at ($(c1 |- regBot)+(0,-0.6)$)
  {\textbf{Generation engine}\\[1pt]\scriptsize several LLMs, including Qwen-2.5; one turn per call, JSON with \texttt{text} and \texttt{intent}\\[1pt] dialogue-state register: goals already met, no repeated questions, product specs in every prompt};
\draw[flow] (rA.south |- regBot) -- (rA.south |- eng.north);
\draw[flow] (rB.south |- regBot) -- (rB.south |- eng.north);

% ======================= phase 2 ====================================================
\node[trj, wide, anchor=north, below=0.35cm of H2] (v1)
  {\textbf{Voice sampling, once per dialogue}\\[1pt]\scriptsize 5 buyer voices and 4 merchant voices, kept fixed for the whole dialogue};
\node[trj, wide, below=of v1] (v2)
  {\textbf{Rewrite with hard rules}\\[1pt]\scriptsize Gemini (\texttt{gemini-3.5-flash-lite}); same turn count and roles; facts, prices, phone numbers and TrxIDs preserved; genuine rewording};
\node[amb, wide, below=of v2] (v3)
  {\textbf{Quality gates}\\[1pt]\scriptsize reject a turn whose SequenceMatcher similarity to its source exceeds 0.82\\[1pt] same amounts and numbers before and after\\[1pt] currency normalised to ``X taka''};
\node[amb, wide, below=of v3] (v4)
  {\textbf{Controlled micro-jitter}\\[1pt]\scriptsize 25\% lowercased first words, 30\% dropped periods, fillers (ok, hmm, well, sure), sparse typos that never touch names or digits};
\draw[flow] (v1) -- (v2); \draw[flow] (v2) -- (v3); \draw[flow] (v3) -- (v4);

% ======================= phase 3 ====================================================
\node[xat, wide, anchor=north, below=0.35cm of H3] (l1)
  {\textbf{Automated semantic labelling}\\[1pt]\scriptsize separate Gemini Flash pass: binary \texttt{buying\_intent}, \texttt{outcome\_category}, a one-sentence \texttt{reason}, normalised turn-level intent tags};
\node[neu, draw=Slate, wide, below=of l1] (l2)
  {\textbf{Human audit of all 700 dialogues}\\[1pt]\scriptsize line-by-line reading; label checked against the buyer's final stance; \texttt{reason} checked for invented facts; borderline cases adjudicated together; \texttt{reversal\_turn\_id} and sentiment tags verified};
\draw[flow] (l1) -- (l2);

% ======================= bottom row: what each phase hands on ======================
% placed below the tallest column (phase 2)
\coordinate (outY) at ($(v4.south)+(0,-0.75)$);
\node[neu, outbox, draw=PhA, anchor=north] (o1) at (c1 |- outY)
  {\textbf{700 draft dialogues}\\[1pt]\scriptsize fluent but rigid: over-grammatical, uniform sentence length, no hedging or typos};
\node[neu, outbox, draw=PhB, anchor=north] (o2) at (c2 |- outY)
  {\textbf{700 humanised dialogues}\\[1pt]\scriptsize informal chat register; facts, labels and turn structure unchanged};
\node[fin, outbox, anchor=north] (o3) at (c3 |- outY)
  {\textbf{Frozen benchmark}\\[1pt]\scriptsize \texttt{ENG\_001}--\texttt{ENG\_500}, \texttt{REV\_001}--\texttt{REV\_200}\\ 6 to 32 turns, 16.2 on average};
\node[neu, fill=white, text width=4.3cm, anchor=north, below=0.3cm of o3] (o4)
  {\scriptsize\textbf{Label counts used in every experiment}\\[1pt] canonical: 262 no / 238 yes\\ reversals: 100 no / 100 yes\\ combined: 362 no / 338 yes};

\draw[flow] (eng) -- (o1);
\draw[flow] (v4) -- (o2);
\draw[flow] (l2) -- (l2.south |- o3.north);
\draw[flow, dsh] (o3) -- (o4);

% ======================= phase frames (fitted, behind) ==============================
\begin{scope}[on background layer]
  \node[pframe=PhA, fit=(H1)(rA)(rB)(eng)(o1)(o4.south -| H1.west)] (F1) {};
  \node[pframe=PhB, fit=(H2)(v1)(v4)(o2)(o4.south -| H2.west)] (F2) {};
  \node[pframe=PhC, fit=(H3)(l1)(l2)(o3)(o4)] (F3) {};
\end{scope}

% ======================= hand-offs between phases ===================================
\draw[handoff] (o1.east) -- (o2.west);
\draw[handoff] (o2.east) -- (o3.west);
\draw[flow, line width=1.3pt, draw=Slate] (H1.east) -- (H2.west);
\draw[flow, line width=1.3pt, draw=Slate] (H2.east) -- (H3.west);

\end{tikzpicture}
}
  \caption{Three-phase synthetic negotiation benchmark generation pipeline, comprising persona-conditioned LLM drafting ($N=500$ canonical, $N=200$ adversarial reversals), organic paraphrasing with linguistic jitter, and multi-tier semantic audit.}
  \label{fig:datagen_pipeline}
\end{figure}

\subsection{The Three-Phase Generation Framework}
\begin{enumerate}
  \item \textbf{Phase 1: Persona-Conditioned Prompt-Driven Synthesis:} Dialogues are generated under two distinct regimes across multiple LLM backbones (including Qwen-2.5):
    \begin{itemize}
      \item \textit{Regime A (Canonical Stream, $N=500$):} Models standard consumer transactions across eight behavioral buyer archetypes (The Skeptic, The LowBaller, The Urgent Buyer, The Ghoster, The Comparator, The DeliveryObsessed, The TrustSeeker, The BulkBuyer), three professional merchant personas (courteous page admin, small shopkeeper, formal retail representative), and four consumer product categories (electronics/phones, traditional apparel, home appliances, organic specialty goods). A Dialogue State Tracking (DST) register prevents repetitive loops.
      \item \textit{Regime B (Adversarial Reversal Stream, $N=200$):} Generates non-monotonic trajectory shifts: 65 positive reversals (\texttt{resistance\_to\_agreement}, where deep skepticism resolves after a verified warranty), 65 negative reversals (\texttt{commit\_to\_rejection}, where early commitment collapses over delivery fees), and 70 hard monotonic controls designed with identical lengths and friction to prevent models from learning length-based shortcuts.
    \end{itemize}
  \item \textbf{Phase 2: Organic Paraphrasing and Humanization:} To strip away synthetic LLM phrasing and replicate mobile messaging dynamics, drafts pass through an organic paraphrasing engine powered by the Gemini API (\texttt{gemini-3.5-flash-lite}). The engine injects controlled linguistic jitter: 25\% lowercased sentence openers, 30\% dropped terminal periods, natural fillers (``hmm'', ``well'', ``ok''), local pricing terms (``X taka''), and sparse typos that strictly protect product names and digits. A strict similarity gate rejects any rewritten turn with a SequenceMatcher ratio $>0.82$ against the source draft.
  \item \textbf{Phase 3: Semantic Labeling and 100\% Team Audit:} An automated pass assigns \texttt{buying\_intent}, multi-class \texttt{outcome\_category}, and textual justifications. The entire research team then conducted an exhaustive 100\% manual audit of all 700 dialogues, verifying stance consistency, correcting subtle hallucinations, and confirming reversal turn IDs ($\tau_{\text{rev}}$).
\end{enumerate}

\subsection{Frozen Benchmark Dataset Counts}
While early design prompts envisioned a 250/250 canonical split across six outcome classes, the finalized, frozen datasets utilized in all experimental runs contain verified counts across four primary outcome categories:
\begin{itemize}
  \item \textbf{Canonical Benchmark ($N=500$):} $262$ \texttt{no} ($52.4\%$) / $238$ \texttt{yes} ($47.6\%$). Outcomes: \texttt{PURCHASE\_COMMITTED} 238, \texttt{DEFERRED\_CONSIDERATION} 139, \texttt{INQUIRY\_DROPOUT} 112, \texttt{EXPLICIT\_REJECTION} 11. Mean turns: $15.9$ (range 6 to 32).
  \item \textbf{Adversarial Reversals ($N=200$):} Exactly $100$ \texttt{no} / $100$ \texttt{yes}. Mean turns: $17.1$ (range 8 to 30).
  \item \textbf{Combined Augmented Benchmark ($N=700$):} $362$ \texttt{no} ($51.7\%$) / $338$ \texttt{yes} ($48.3\%$). Outcomes: \texttt{PURCHASE\_COMMITTED} 338, \texttt{DEFERRED\_CONSIDERATION} 157, \texttt{INQUIRY\_DROPOUT} 125, \texttt{EXPLICIT\_REJECTION} 80. Mean turns: $16.2$ (range 6 to 32).
\end{itemize}

\subsection{Generic Architecture Diagnostic and Ablation Harness}
\label{sec:diagnostic-ablation-harness}

To determine whether the architectural mechanisms in Antiphony (speaker disentanglement, directional cross-attention, and trajectory tracking) reflect genuine domain-agnostic inductive biases for asymmetric multi-turn dialogue rather than overfitting to donation persuasion, we engineered a dedicated generic architecture diagnostic and ablation harness.

\subsubsection{Counterfactual Input-Ablation Operators}
In conversational benchmarks, conversational outcomes can easily become overdetermined: explicit affirmations or closing logistical exchanges allow flat sequence classifiers to infer intent from isolated surface tokens, bypassing dialogue history. To systematically pinpoint the causal source of predictive representations, the harness formalizes four distinct counterfactual input-ablation operators:
\begin{enumerate}
  \item \textbf{Last-Turn-Removed ($\text{Ablation}_{\text{LTR}}$):} Excises the final decider turn:
    \begin{equation}
      \mathcal{D}_{\text{LTR}} = \mathcal{D} \setminus \{u_{\text{final}}\}
    \end{equation}
    This evaluates whether preceding conversational momentum enables outcome inference without an explicit closing declaration.
  \item \textbf{Last-Turn and Merchant Courtesy Removed ($\text{Ablation}_{\text{LTR+M}}$):} In commercial transactions, merchants routinely emit closing courtesy affirmations (payment receipts, shipping notices) immediately following agreement. This operator excises both the terminal buyer turn and any trailing merchant courtesy turns:
    \begin{equation}
      \mathcal{D}_{\text{LTR+M}} = \mathcal{D} \setminus \{u \in \mathcal{D} \mid \text{index}(u) \ge \text{index}(u_{\text{final}})\}
    \end{equation}
    thereby completely severing closing bidirectional lexical leakage.
  \item \textbf{First-Half-Only ($\text{Ablation}_{\text{FHO}}$):} Dialogues are restricted strictly to the first $50\%$ of turns ($u_1, \dots, u_{\lfloor K/2 \rfloor}$), detecting early predictive bias and verifying that synthetic personas do not leak the outcome in the initial exchanges.
  \item \textbf{Pre-Reversal-Only ($\text{Ablation}_{\text{PRO}}$):} Applied to adversarial reversal dialogues exhibiting an annotated stance-shift at turn $\tau_{\text{rev}}$, this operator excises all content from the shift onward:
    \begin{equation}
      \mathcal{D}_{\text{PRO}} = [u_1, \dots, u_{\tau_{\text{rev}} - 1}]
    \end{equation}
    testing whether models commit prematurely to initial buyer resistance or correctly exploit merchant accommodation signals.
\end{enumerate}

\subsubsection{Tokenizer Truncation Directionality Invariant}
During the engineering of the ablation harness, a critical diagnostic defect was exposed: standard transformer tokenization libraries apply right-truncation by default (\texttt{truncation\_side='right'}, capped at 512 tokens). In lengthy multi-turn commercial chats ($30.4\%$ of dialogues in the benchmark exceeded 512 tokens), right-truncation silently discards the terminal tail of the conversation. When applying $\text{Ablation}_{\text{LTR}}$, right-truncated tokenizers mistakenly pruned the proximal turns preceding the excised turn, generating an artificial $+0.1600$ necessity collapse (Macro-F1 plunging to $0.8133$).

The diagnostic harness enforces a strict \textbf{left-truncation invariant} (\texttt{truncation\_side='left'}):
\begin{equation}
  \text{Tokens}(\mathcal{D}) = [t_{\max(1, N_{\text{tokens}} - L_{\max} + 1)}, \dots, t_{N_{\text{tokens}}}]
\end{equation}
which discards only early introductory pleasantries while strictly preserving the vital decision-boundary turns. Enforcing this invariant recovered baseline performance to $0.9331$, establishing that the true necessity gap on canonical dialogues is modest ($+0.0402$).

\subsubsection{Parametric Skew-Stress Formulation}
To evaluate operational resilience against severe real-world class imbalance without retraining the shared encoder, the harness implements a parametric class-pruning generator. Given a balanced training distribution ($47.7\%$ affirmative, $n=350$), affirmative dialogues are subsampled according to target ratio $\alpha \in [0.10, 0.50]$:
\begin{equation}
  n_{\text{pos}} = \left\lfloor \frac{\alpha}{1 - \alpha} \cdot n_{\text{neg}} \right\rfloor
\end{equation}
depressing the affirmative training ratio down to $9.85\%$ ($n=203$). This profiles how gracefully directional cross-attention and recurrent trajectory tracking maintain calibration compared to flat sequence models under operational distribution collapse.

\section{Strategy and Category Modeling Architecture}
\label{sec:strategy-modeling}

Predicting multi-label rhetorical strategies across an extreme long tail requires a multi-task formulation supported by an ensemble of diverse base learners and a calibrated decision layer.

\begin{figure}[htbp]
  \centering
  \resizebox{\textwidth}{!}{% Strategy / category classifier (v18): how the 19-member pool is built and decided.
% Three stacked bands, read top to bottom:
%   1. one ensemble member: the transformer recipe (runs 7 times)
%   2. the pool of 19 members, grouped by how each is built
%   3. the decision layer that turns 19 probabilities into a label set
% Uses the shared palette and styles from tikzstyles.tex (amber = strategy stream,
% green = few-shot companions, neutral = classical / zero-shot, slate = outputs).
% Width is about 16.4 cm; needs amssymb for \mathbb.
\begin{tikzpicture}[x=1cm,y=1cm,
  band/.style={rounded corners=6pt, line width=0.7pt},
  btitle/.style={font=\footnotesize\bfseries, inner sep=2pt},
  sub/.style={font=\scriptsize},
  rstep/.style={amb, text width=1.5cm, minimum height=1.55cm, font=\scriptsize},
  member/.style={minimum height=2.1cm},
  bigflow/.style={flow, line width=1.1pt},
]

% ======================= band backdrops =============================================
\draw[band, draw=AmbBord!60, fill=AmbFill!45] (3.85,-1.0) rectangle (12.95,-4.45);
\draw[band, draw=NeuBord!60, fill=NeuFill!60] (0,-5.35) rectangle (16.4,-8.3);
\draw[band, draw=Slate!55, fill=NeuFill!35] (0,-9.85) rectangle (16.4,-13.15);

\node[btitle, text=AmbBord, anchor=north east] at (12.8,-1.08)
  {One member: the transformer recipe\quad{\normalfont\scriptsize runs 7 times}};
\node[btitle, text=Slate, anchor=north] at (9.2,-5.42) {Pool of 19 members};
\node[btitle, text=Slate, anchor=north west] at (4.55,-9.92) {Decision layer};

% ======================= input ======================================================
\node[neu, text width=3.1cm] (in) at (1.85,0)
  {\textbf{Input}\\[1pt]\scriptsize a Persuader turn plus up to five turns of context};

% ======================= band 1: transformer recipe =================================
\node[rstep] (t1) at (4.99,-2.6)
  {\textbf{Pretrained transformer}\\[1pt] RoBERTa-base, RoBERTa-large or ELECTRA-base};
\node[rstep] (t2) at (7.13,-2.6)
  {\textbf{Role embedding}\\[1pt] added to every token: Persuader or Persuadee};
\node[rstep] (t3) at (9.27,-2.6)
  {\textbf{Attention pooling}\\[1pt] one vector $z\in\mathbb{R}^d$, $d{=}768$ or $1024$};
\node[amb, text width=1.9cm, font=\scriptsize] (hc) at (11.61,-2.02)
  {\textbf{Category head}\\ \mbox{\texttt{Linear}($d$,\,11)}\\ main target};
\node[amb, fill=white, text width=1.9cm, font=\scriptsize] (hs) at (11.61,-3.2)
  {\textbf{Strategy head}\\ \mbox{\texttt{Linear}($d$,\,41)}\\ auxiliary};
\draw[flow] (t1) -- (t2);
\draw[flow] (t2) -- (t3);
\draw[flow] (t3.east) -- ++(0.18,0) |- (hc.west);
\draw[flow] (t3.east) -- ++(0.18,0) |- (hs.west);
\node[nt, anchor=south] at (8.4,-4.42)
  {asymmetric loss on both heads; rare-label copies in the training split only};

% ======================= TF-IDF vectorizer (shared features, not a member) ==========
\node[neu, fill=white, dsh, text width=2.6cm] (vec) at (14.72,-2.6)
  {\textbf{TF-IDF vectorizer}\\[1pt]\scriptsize word and character n-grams; shared features, not a member};

% ======================= band 2: the 19 members =====================================
\node[neu, fill=white, member, text width=3.0cm] (zs) at (1.85,-6.95)
  {\textbf{2 zero-shot members}\\[1pt]\scriptsize similarity between the turn and each label's written definition\\[1pt]\textit{no training at all}};
\node[amb, member, text width=3.75cm] (enc) at (5.9,-6.95)
  {\textbf{7 fine-tuned encoders}\\[1pt]\scriptsize RoBERTa-base $\times4$, RoBERTa-large $\times2$, ELECTRA-base $\times1$; mixed context width, seed and loss setting\\[1pt]\textit{the only members that run the recipe}};
\node[trj, member, text width=3.95cm] (fs) at (10.75,-6.95)
  {\textbf{8 few-shot companions}\\[1pt]\scriptsize distance to each label's average training vector; 7 on a frozen encoder vector $z$, 1 on TF-IDF\\[1pt]\textit{no gradient step of their own}};
\node[neu, member, text width=2.6cm] (tc) at (14.72,-6.95)
  {\textbf{2 TF-IDF classifiers}\\[1pt]\scriptsize logistic regression and XGBoost\\[1pt]\textit{trained from scratch}};

% ======================= band 3: decision layer =====================================
\node[fin, text width=15.5cm] (probs) at (8.2,-9.1)
  {19 probabilities for every (turn, category) pair\quad{\scriptsize\itshape how a member was built no longer matters from here}};

\node[neu, fill=white, text width=3.3cm, minimum height=2.3cm] (ens) at (2.15,-11.55)
  {\textbf{Ensemble search}\\[1pt]\scriptsize 50 to 70 candidate combiners (averages, greedy and bagged blends, power means, two stackers), scored by grouped cross-validation};
\node[amb, fill=white, line width=1.3pt, text width=2.6cm] (bl) at (6.4,-10.95)
  {\textbf{Calibrated blend}\\\scriptsize $q{=}2$ power mean\\ \textbf{selected}\,$\cdot$\,0.572};
\node[amb, fill=white, dsh, text width=2.6cm] (xg) at (6.4,-12.35)
  {\textbf{XGBoost stacker}\\\scriptsize tuned, logits only\\ tested\,$\cdot$\,0.551};
\node[neu, fill=white, text width=2.9cm] (rule) at (10.5,-10.95)
  {\textbf{Per-label decision rule}\\[1pt]\scriptsize chosen from nine candidates on validation data};
\node[fin, text width=3.0cm] (out) at (14.55,-10.95)
  {\textbf{Label set for the turn}\\[1pt]\scriptsize over the 11 categories};
\node[nt, anchor=west, align=left] at (7.95,-12.35) {scores are cross-fitted\\validation macro-F1};

% ======================= arrows =====================================================
% input feeds the recipe, the TF-IDF vectorizer and the zero-shot members
\draw[flow] (in.east) -- (14.72,0) -- (vec.north);
\draw[flow] (4.99,0) -- (t1.north);
\fill[SlateSoft] (4.99,0) circle (1.4pt);
\draw[flow] (in.south) -- (zs.north);

% recipe -> the 7 encoders
\draw[bigflow, draw=AmbMain] (5.9,-4.45) -- node[right, font=\scriptsize\bfseries, text=AmbBord]{$\times 7$} (enc.north);

% TF-IDF vectorizer -> TF-IDF classifiers and the 8th few-shot companion
\draw[flow] (vec.south) -- (tc.north);
\draw[flow] (14.72,-4.9) -- (12.3,-4.9) -- (12.3,-5.85);
\fill[SlateSoft] (14.72,-4.9) circle (1.4pt);

% encoders -> few-shot companions (frozen vector z)
\draw[flowT] (enc.east) -- node[above, font=\tiny, text=TrjBord]{$z$} (fs.west);

% every member group -> the 19 probabilities
\foreach \m in {zs,enc,fs,tc}{ \draw[flow] (\m.south) -- (\m.south |- probs.north); }

% decision layer
\draw[flow] (probs.south -| ens.north) -- (ens.north);
\draw[flow] (ens.east) -- ++(0.2,0) |- (bl.west);
\draw[flow] (ens.east) -- ++(0.2,0) |- (xg.west);
\draw[bigflow, draw=AmbMain] (bl) -- (rule);
\draw[bigflow] (rule) -- (out);

\end{tikzpicture}
}
  \caption{Strategy and category classification framework, illustrating the 19-member model pool, multi-task representation heads, and decision layer search comparing calibrated power-mean blending (selected, val macro-F1 0.5716) against XGBoost stacking (tested, val macro-F1 0.5508).}
  \label{fig:strategy_model}
\end{figure}

\subsection{Multi-Task Category-First Formulation}
Because the 41 strategies include classes with as few as 2 examples (\texttt{door\_in\_the\_face}), models trained solely on 41 fine strategies achieve poor macro-F1 ($0.368$). We pivot the primary modeling target to the 11 communicative categories, where every category has over 100 real instances. We formulate a multi-task network where the 11 categories act as the primary objective ($\mathcal{L}_{\text{cat}}$) and the 41 strategies act as an auxiliary regularizer ($\mathcal{L}_{\text{strat}}$):
\begin{equation}
  \mathcal{L}_{\text{total}} = \mathcal{L}_{\text{cat}} + \lambda_{\text{aux}} \mathcal{L}_{\text{strat}}
\end{equation}
with $\lambda_{\text{aux}} \in [0.1, 0.3]$. Both heads are optimized using Asymmetric Loss (ASL) \cite{benbaruch2021asymmetric}:
\begin{equation}
  \mathcal{L}_{\text{ASL}} = -\sum_{c=1}^{C} \left[ y_c (1 - p_c)^{\gamma_+} \log(p_c) + (1 - y_c) (p_{m,c})^{\gamma_-} \log(1 - p_{m,c}) \right]
\end{equation}
where $p_{m,c} = \max(p_c - m, 0)$ incorporates hard probability margin shifting ($m=0.05$) and asymmetric focusing ($\gamma_+ = 0, \gamma_- = 4$) to discard easy negative gradients.

\subsection{The 19-Member Heterogeneous Model Pool}
To maximize ensemble diversity, we construct a 19-member candidate pool:
\begin{itemize}
  \item \textbf{7 Fine-Tuned Transformer Encoders:} RoBERTa-base (4 variations with mixed random seeds, context widths, and loss weighting), RoBERTa-large (2 models, with and without dialogue context), and ELECTRA-base (1 model).
  \item \textbf{8 Few-Shot Vector Companions:} Nearest-centroid cosine distance classifiers over mean-pooled encoder representations (one per transformer, plus one over TF-IDF space).
  \item \textbf{2 Classical TF-IDF Baselines:} Character- and word-level n-gram TF-IDF models paired with Logistic Regression and XGBoost.
  \item \textbf{2 Zero-Shot Semantic Similarity Models:} Lexical cue matchers mapping turn embeddings to category anchor definitions.
\end{itemize}

\subsection{Decision Layer Search: Calibrated Blend vs. XGBoost Stacker}
The predictions of the 19 base members were combined by evaluating a wide zoo of decision combiners across grouped 5-fold cross-validation. A central research question was whether a machine-learned XGBoost stacking layer could replace standard probability blending:
\begin{itemize}
  \item \textbf{XGBoost Stacking Experiment (Tested, Not Selected):} We constructed a tabular stacking dataset of $\approx 16{,}600$ (turn, label) instances across two feature regimes: `\texttt{logits mode'' (20 raw member pre-sigmoid outputs) and }`rich mode'' (41 features including member disagreement variances, parent category scores, and intra-turn ranks). Hyperparameter grid search over 72 configurations (max depth 4, learning rate 0.06, 300 to 500 trees) achieved a tuned cross-validated Macro-F1 of \textbf{$0.5508$} (logits mode) and micro-F1 of \textbf{$0.6847$} (rich mode), significantly improving over the untuned baseline ($0.5329$).
  \item \textbf{Calibrated Power-Mean Blending (Selected Winner):} Despite extensive tuning, the XGBoost stacker was outperformed by a \textbf{calibrated power-mean blend} ($q=+2.0$) applied to member probabilities post-Platt scaling \cite{platt1999probabilistic}:
    \begin{equation}
      \bar{p}_c = \left( \frac{1}{M} \sum_{m=1}^{M} \left[ \sigma(\alpha_m z_{m,c} + \beta_m) \right]^q \right)^{1/q}
    \end{equation}
    The calibrated power-mean blend achieved a superior cross-validated Macro-F1 of \textbf{$0.5716$}. Consequently, the calibrated blend was selected as the shipped decision layer.
\end{itemize}

\clearpage
\section{The AntiphonyClassifier Architecture}
\label{sec:antiphony}

The centerpiece intention architecture of this project is the \textit{AntiphonyClassifier}. Built to overcome the limitations of flat sequence transformers and symmetric dialogue encoders, Antiphony explicitly embeds the fundamental asymmetry of persuasion into its neural wiring.

To provide both clear intuitive comprehension and implementation-grade precision, we present Antiphony across two complementary architectural schematics:
\begin{itemize}
  \item \textbf{Conceptual Architecture (\Cref{fig:antiphony_simple}):} Details the intuitive functional dataflow—from alternating raw dialogue turns and role disentanglement to directional question-argument cross-attention and dual temporal-pooling context readouts.
  \item \textbf{Tensor and Layer Architecture (\Cref{fig:antiphony_arch}):} Specifies exact mathematical dimensions, transformer sub-layer chains, multi-head projection matrices, residual normalizations, and composite loss optimization heads.
\end{itemize}

\begin{figure}[htbp]
  \centering
  \resizebox{\textwidth}{!}{% AntiphonyClassifier (v14), conceptual view: the same blocks as antiphony_arch.tikz,
% described in words instead of equations. Same palette, same top-to-bottom order.
% Drawn at 16 cm width.
\begin{tikzpicture}[x=1cm,y=1cm,
  why/.style={font=\scriptsize\itshape},
  lbl/.style={font=\scriptsize\itshape, text=SlateSoft, fill=white, inner sep=1.5pt},
  chip/.style={minimum width=1.55cm, minimum height=0.46cm, inner sep=1pt, font=\scriptsize},
]

% ---------------- input dialogue ----------------------------------------------
\node[neu, minimum width=15.6cm, minimum height=1.35cm] (dlg) at (8,0) {};
\node[ttl] at (8,0.36) {Input: one dialogue, turns alternating between the two speakers};
\foreach \i/\x in {1/1.6,3/5.2,5/8.8}{
  \node[per, chip] at (\x,-0.25) {Persuader};}
\foreach \i/\x in {2/3.4,4/7.0,6/10.6}{
  \node[dec, chip] at (\x,-0.25) {Persuadee};}
\node[font=\scriptsize] at (12.0,-0.25) {$\cdots$};
\node[dec, chip] at (13.4,-0.25) {last turn};

% ---------------- turn encoder and split ------------------------------------------
\node[neu, minimum width=15.6cm, minimum height=0.9cm] (enc) at (8,-1.75)
  {\textbf{Turn encoder}\quad a pretrained language model (RoBERTa, DeBERTa-v3 or TOD-BERT) reads each turn and summarises it as one vector};
\draw[flow] (dlg) -- (enc);
\node[neu, minimum width=6.2cm] (split) at (8,-2.95) {Separate the turns by speaker, keeping their order};
\draw[flow] (enc) -- (split);

% ---------------- two speaker streams ---------------------------------------------
\node[per, text width=6.9cm, minimum height=1.55cm] (sP) at (4.0,-4.55)
  {\textbf{Persuader stream}\\[1pt]\scriptsize a small Transformer reads only the Persuader's turns, in order\\[2pt] {\scriptsize\itshape\color{SlateSoft} how does the argument develop?}};
\node[dec, text width=6.9cm, minimum height=1.55cm] (sD) at (12.0,-4.55)
  {\textbf{Persuadee stream}\\[1pt]\scriptsize a small Transformer reads only the Persuadee's turns, in order\\[2pt] {\scriptsize\itshape\color{SlateSoft} how does the decision-maker's stance develop?}};
\draw[flowP] (split.west) -| (sP.north);
\draw[flowD] (split.east) -| (sD.north);

% ---------------- directional cross-attention -------------------------------------
\node[xat, text width=15.1cm, minimum height=1.6cm] (xa) at (8,-6.95)
  {\textbf{Directional cross-attention}\quad{\scriptsize the Persuadee asks, the Persuader answers}\\[2pt]
   \scriptsize each Persuadee turn looks back over the Persuader's turns and pulls in the arguments it is responding to;\\
   \scriptsize the result is added back to the Persuadee's own representation\\[2pt]
   {\scriptsize\itshape\color{SlateSoft} given what was argued, how is the decision-maker responding?}};
\draw[flowP] (sP.south) -- node[lbl, right] {arguments} (sP.south |- xa.north);
\draw[flowD] (sD.south) -- node[lbl, left] {questions} (sD.south |- xa.north);

\node[ten, fill=XatFill, font=\scriptsize] (ctx) at (8,-8.45) {Persuadee turns, now aware of what was argued};
\draw[flowX] (xa) -- (ctx);

% ---------------- two readouts --------------------------------------------------
\node[trj, text width=6.9cm, minimum height=1.75cm] (gru) at (4.0,-10.0)
  {\textbf{Commitment trajectory}\enspace{\scriptsize(GRU)}\\[1pt]\scriptsize reads the Persuadee's turns in order with a running memory, and keeps that memory as it stands at the last Persuadee turn\\[2pt] {\scriptsize\itshape\color{SlateSoft} how did commitment build up, turn by turn?}};
\node[trj, text width=6.9cm, minimum height=1.75cm] (pool) at (12.0,-10.0)
  {\textbf{Attention pooling}\\[1pt]\scriptsize gives every Persuadee turn an importance weight and takes the weighted average\\[2pt] {\scriptsize\itshape\color{SlateSoft} which moments mattered most overall?}};
\draw[flowX] (ctx.west) -| (gru.north);
\draw[flowX] (ctx.east) -| (pool.north);

% ---------------- fusion, head, output --------------------------------------------
\node[fin, minimum width=5.4cm] (cat) at (8,-11.65) {Combine both views};
\draw[flowT] (gru.south) |- (cat.west);
\draw[flowT] (pool.south) |- (cat.east);
\node[fin, minimum width=5.4cm] (head) at (8,-12.65) {Classifier (2-layer MLP)};
\node[fin, minimum width=5.4cm] (yh) at (8,-13.65) {Does the Persuadee commit to donate?\enspace yes / no};
\draw[flow] (cat) -- (head); \draw[flow] (head) -- (yh);

% ---------------- legend ------------------------------------------------------------
\begin{scope}[shift={(0.5,-14.75)}]
  \node[per, minimum width=0.45cm, minimum height=0.3cm] at (0,0) {};
  \node[anchor=west, font=\scriptsize] at (0.35,0) {Persuader};
  \node[dec, minimum width=0.45cm, minimum height=0.3cm] at (2.4,0) {};
  \node[anchor=west, font=\scriptsize] at (2.75,0) {Persuadee};
  \node[xat, minimum width=0.45cm, minimum height=0.3cm] at (4.9,0) {};
  \node[anchor=west, font=\scriptsize] at (5.25,0) {Cross-attention};
  \node[trj, minimum width=0.45cm, minimum height=0.3cm] at (8.1,0) {};
  \node[anchor=west, font=\scriptsize] at (8.45,0) {Trajectory / pooling};
  \node[fin, minimum width=0.45cm, minimum height=0.3cm] at (11.7,0) {};
  \node[anchor=west, font=\scriptsize] at (12.05,0) {Fusion / output};
\end{scope}

\end{tikzpicture}
}
  \caption{High-level conceptual schematic of the AntiphonyClassifier architecture, illustrating the intuitive information flow from alternating conversational turns to speaker stream disentanglement, directional asymmetric cross-attention, dual trajectory and attention pooling context aggregation, and final commitment classification.}
  \label{fig:antiphony_simple}
\end{figure}

\subsection{Detailed Mathematical Formulation}
Let dialogue $\mathcal{D}$ comprise $U$ utterances ($U \le 32$), each truncated to $L=64$ tokens.
\begin{enumerate}
  \item \textbf{Batched Utterance Encoding:} The input tensor $\mathbf{X} \in \mathbb{R}^{B \times U \times L}$ is flattened to $[B \cdot U, L]$ and passed through a single batched forward pass of a shared transformer encoder (e.g., DeBERTa-v3-base), yielding token hidden states $\mathbf{E} \in \mathbb{R}^{(B \cdot U) \times L \times d}$ where $d=768$. Masked mean-pooling over valid non-padded tokens produces utterance vectors $\mathbf{U} \in \mathbb{R}^{B \times U \times d}$.
  \item \textbf{Role-Pure Disentanglement:} Utterances are partitioned using binary role masks into two role-pure streams: Persuader turns $\mathbf{H}^{(\mathcal{P})} \in \mathbb{R}^{B \times R_{\mathcal{P}} \times d}$ and Persuadee/Decider turns $\mathbf{H}^{(\mathcal{D})} \in \mathbb{R}^{B \times R_{\mathcal{D}} \times d}$, where each stream retains the last $R=20$ turns in chronological order with stream-local position embeddings.
  \item \textbf{Independent Stream Transformers:} Each stream is encoded by an independent 2-layer transformer with multi-head self-attention, pre-LayerNorm, and feed-forward sub-layers:
    \begin{equation}
      \widetilde{\mathbf{H}}^{(\mathcal{P})} = \text{Transformer}_{\mathcal{P}}(\mathbf{H}^{(\mathcal{P})}), \quad \widetilde{\mathbf{H}}^{(\mathcal{D})} = \text{Transformer}_{\mathcal{D}}(\mathbf{H}^{(\mathcal{D})})
    \end{equation}
  \item \textbf{Directional Asymmetric Cross-Attention:} To model how decider stance evolves in response to persuader arguments, the Decider stream acts as the Query, while the Persuader stream provides the Keys and Values:
    \begin{equation}
      \mathbf{Q} = \widetilde{\mathbf{H}}^{(\mathcal{D})} \mathbf{W}_Q, \quad \mathbf{K} = \widetilde{\mathbf{H}}^{(\mathcal{P})} \mathbf{W}_K, \quad \mathbf{V} = \widetilde{\mathbf{H}}^{(\mathcal{P})} \mathbf{W}_V
    \end{equation}
    \begin{equation}
      \mathbf{A}^{(\mathcal{D} \to \mathcal{P})} = \text{softmax}\left(\frac{\mathbf{Q} \mathbf{K}^\top}{\sqrt{d_k}} + \mathbf{M}_{\mathcal{P}}\right), \quad \mathbf{C}^{(\mathcal{D} \to \mathcal{P})} = \mathbf{A}^{(\mathcal{D} \to \mathcal{P})} \mathbf{V}
    \end{equation}
    The cross-attended representations are combined via residual layer normalization:
    \begin{equation}
      \mathbf{X} = \text{LayerNorm}(\widetilde{\mathbf{H}}^{(\mathcal{D})} + \mathbf{C}^{(\mathcal{D} \to \mathcal{P})})
    \end{equation}
  \item \textbf{Dual Context Aggregation:} The contextualized decider representations $\mathbf{X} = [\mathbf{x}_1, \dots, \mathbf{x}_T]$ are aggregated across two complementary paths:
    \begin{itemize}
      \item \textit{Commitment-Trajectory GRU:} Captures non-monotonic stance shifts via a recurrent hidden state:
        \begin{equation}
          \mathbf{h}_t = \text{GRU}(\mathbf{x}_t, \mathbf{h}_{t-1})
        \end{equation}
        The trajectory state $\mathbf{h}_T$ is extracted at the dialogue's true last valid decider turn, bypassing trailing padding slots.
      \item \textit{Learned Attention Pooling:} Computes an order-invariant summary of salient turns:
        \begin{equation}
          \alpha_t = \frac{\exp(\mathbf{w}_a^\top \tanh(\mathbf{W}_a \mathbf{x}_t))}{\sum_{j=1}^{T} \exp(\mathbf{w}_a^\top \tanh(\mathbf{W}_a \mathbf{x}_j))}, \quad \mathbf{p} = \sum_{t=1}^{T} \alpha_t \mathbf{x}_t
        \end{equation}
    \end{itemize}
  \item \textbf{Fusion and Classification Head:} The temporal trajectory vector $\mathbf{h}_T$ and the attention-pooled vector $\mathbf{p}$ are concatenated:
    \begin{equation}
      \mathbf{z} = [\mathbf{h}_T \,;\, \mathbf{p}] \in \mathbb{R}^{2d}
    \end{equation}
    and fed into a 2-layer MLP head ($\text{Linear}(2d, d) \to \text{GELU} \to \text{Dropout}(0.2) \to \text{Linear}(d, 2)$) producing unnormalized binary logits $\hat{\mathbf{y}} \in \mathbb{R}^2$.
  \item \textbf{Multi-Objective Training Loss:} Antiphony is optimized using a joint composite loss combining Class-Balanced Focal Loss \cite{cui2019class} ($\gamma=2.0, \beta=0.999$) and Supervised Contrastive Loss \cite{khosla2020supervised} ($\tau=0.1$, weight $\lambda_{\text{sc}} = 0.30$):
    \begin{equation}
      \mathcal{L} = \mathcal{L}_{\text{CB-Focal}}(\hat{\mathbf{y}}, y) + \lambda_{\text{sc}} \mathcal{L}_{\text{SupCon}}(\mathbf{z}, y)
    \end{equation}
\end{enumerate}

\begin{figure}[p]
  \centering
  \resizebox{\textwidth}{!}{% AntiphonyClassifier (v14): role-pure streams, directional cross-attention,
% commitment-trajectory GRU and parallel attention pooling.
% Drawn at 16 cm width; tensor names follow Section 4.6 of the report.
% Needs amsmath/amssymb (\mathbb, \operatorname) in the preamble.
\begin{tikzpicture}[x=1cm,y=1cm]

% ---------------- group backdrops (drawn first, nodes sit on top) --------------
\draw[rounded corners=6pt, draw=PerBord!60, fill=PerFill!40, line width=0.7pt] (0.2,-4.35) rectangle (7.6,-7.55);
\draw[rounded corners=6pt, draw=DecBord!60, fill=DecFill!40, line width=0.7pt] (8.4,-4.35) rectangle (15.8,-7.55);
\draw[rounded corners=6pt, draw=XatBord!70, fill=XatFill!45, line width=0.8pt] (0.2,-8.25) rectangle (15.8,-13.15);
\draw[rounded corners=6pt, draw=TrjBord!60, fill=TrjFill!40, line width=0.7pt] (0.2,-13.85) rectangle (7.6,-17.0);
\draw[rounded corners=6pt, draw=TrjBord!60, fill=TrjFill!40, line width=0.7pt] (8.4,-13.85) rectangle (15.8,-17.0);

% ---------------- input dialogue ----------------------------------------------
\node[neu, minimum width=15.6cm, minimum height=1.5cm] (dlg) at (8,0) {};
\node[ttl] at (8,0.42) {Dialogue $\mathcal{D}=(u_1,\dots,u_N)$, at most $32$ turns of $64$ tokens, role $r_i\in\{P,D\}$};
\foreach \i [evaluate=\xx using {1.4+1.5*\i}] in {1,3,5}{
  \node[per, minimum width=1.1cm, minimum height=0.48cm, inner sep=1pt, font=\scriptsize] at (\xx,-0.28) {$u_{\i}$};}
\foreach \i [evaluate=\xx using {1.4+1.5*\i}] in {2,4,6}{
  \node[dec, minimum width=1.1cm, minimum height=0.48cm, inner sep=1pt, font=\scriptsize] at (\xx,-0.28) {$u_{\i}$};}
\node[font=\scriptsize] at (11.7,-0.28) {$\cdots$};
\node[dec, minimum width=1.1cm, minimum height=0.48cm, inner sep=1pt, font=\scriptsize] at (13.0,-0.28) {$u_N$};

% ---------------- shared utterance encoder ---------------------------------------
\node[neu, minimum width=15.6cm, minimum height=1.0cm, align=center] (enc) at (8,-2.0)
  {\textbf{Shared utterance encoder} (RoBERTa-base $\mid$ DeBERTa-v3-base $\mid$ TOD-BERT): one batched pass over $[B{\cdot}U,\,L]$, masked mean-pool per turn\\[-1pt]
   {\scriptsize $\longrightarrow\ \mathbf{U}\in\mathbb{R}^{B\times U\times d}$}};
\draw[flow] (dlg.south) -- (enc.north);

% ---------------- role split ------------------------------------------------------
\node[neu, minimum width=6.2cm] (split) at (8,-3.2) {Role split: last $R{=}20$ turns per role, chronological};
\draw[flow] (enc.south) -- (split.north);
\node[ten] (hP) at (4.3,-4.0) {$\mathbf{H}^{(P)}\in\mathbb{R}^{R\times d}$};
\node[ten] (hD) at (12.5,-4.0) {$\mathbf{H}^{(D)}\in\mathbb{R}^{R\times d}$};
\draw[flowP] (split.west) -| (hP.north);
\draw[flowD] (split.east) -| (hD.north);

% ---------------- stream transformers ---------------------------------------------
% titles sit top-left in each band, clear of the vertical tensor path
\node[ttl, text=PerBord, anchor=north west, align=left] at (0.35,-4.45) {Persuader stream\\Transformer ($\times 2$)};
\node[ttl, text=DecBord, anchor=north west, align=left] at (8.55,-4.45) {Persuadee stream\\Transformer ($\times 2$)};
\node[ten] (pP) at (4.3,-5.3) {$\mathbf{H}^{(P)}+\mathbf{E}^{(P)}_{\mathrm{pos}}$};
\node[ten] (pD) at (12.5,-5.3) {$\mathbf{H}^{(D)}+\mathbf{E}^{(D)}_{\mathrm{pos}}$};
\draw[flowP] (hP.south) -- (pP.north);
\draw[flowD] (hD.south) -- (pD.north);
% sub-layer chain, persuader
\node[per, minimum width=1.55cm, minimum height=0.85cm, font=\scriptsize] (aP1) at (1.25,-6.3) {Multi-head\\self-attn.};
\node[per, minimum width=1.35cm, minimum height=0.85cm, font=\scriptsize] (aP2) at (3.15,-6.3) {Add \&\\LayerNorm};
\node[per, minimum width=1.35cm, minimum height=0.85cm, font=\scriptsize] (aP3) at (4.85,-6.3) {Feed-\\forward};
\node[per, minimum width=1.35cm, minimum height=0.85cm, font=\scriptsize] (aP4) at (6.55,-6.3) {Add \&\\LayerNorm};
\draw[flowP] (pP.west) -| (aP1.north);
\draw[flowP] (aP1) -- (aP2); \draw[flowP] (aP2) -- (aP3); \draw[flowP] (aP3) -- (aP4);
% sub-layer chain, persuadee
\node[dec, minimum width=1.55cm, minimum height=0.85cm, font=\scriptsize] (aD1) at (9.45,-6.3) {Multi-head\\self-attn.};
\node[dec, minimum width=1.35cm, minimum height=0.85cm, font=\scriptsize] (aD2) at (11.35,-6.3) {Add \&\\LayerNorm};
\node[dec, minimum width=1.35cm, minimum height=0.85cm, font=\scriptsize] (aD3) at (13.05,-6.3) {Feed-\\forward};
\node[dec, minimum width=1.35cm, minimum height=0.85cm, font=\scriptsize] (aD4) at (14.75,-6.3) {Add \&\\LayerNorm};
\draw[flowD] (pD.west) -| (aD1.north);
\draw[flowD] (aD1) -- (aD2); \draw[flowD] (aD2) -- (aD3); \draw[flowD] (aD3) -- (aD4);
% outputs of the streams
\node[ten] (tP) at (4.3,-7.15) {$\tilde{\mathbf{H}}^{(P)}\in\mathbb{R}^{R\times d}$};
\node[ten] (tD) at (12.5,-7.15) {$\tilde{\mathbf{H}}^{(D)}\in\mathbb{R}^{R\times d}$};
\draw[flowP] (aP4.south) |- (tP.east);
\draw[flowD] (aD4.south) |- (tD.east);

% ---------------- cross-attention block -------------------------------------------
\node[ttl, text=XatBord, anchor=south west, align=left] at (0.4,-13.05) {Directional asymmetric cross-attention\\[-1pt]{\scriptsize\normalfont the Persuadee queries the Persuader ($4$ heads)}};
\node[per, minimum width=2.9cm] (wV) at (1.9,-9.3)
  {{\scriptsize\bfseries Value}\enspace$\mathbf{V}=\tilde{\mathbf{H}}^{(P)}\mathbf{W}_V$};
\node[per, minimum width=2.9cm] (wK) at (5.0,-9.3)
  {{\scriptsize\bfseries Key}\enspace$\mathbf{K}=\tilde{\mathbf{H}}^{(P)}\mathbf{W}_K$};
\node[dec, minimum width=2.9cm] (wQ) at (12.5,-9.3)
  {{\scriptsize\bfseries Query}\enspace$\mathbf{Q}=\tilde{\mathbf{H}}^{(D)}\mathbf{W}_Q$};
\draw[flowP] (tP.south) -- ++(0,-0.4) -| (wK.north);
\draw[flowP, dsh] (tP.south) -- ++(0,-0.4) -| (wV.north);
\draw[flowD] (tD.south) -- (wQ.north);

\node[xat, minimum width=4.9cm] (att) at (8.6,-10.45)
  {$\mathbf{A}=\operatorname{softmax}\!\big(\mathbf{Q}\mathbf{K}^{\top}/\sqrt{d_k}\big)\in\mathbb{R}^{R_D\times R_P}$};
\draw[flowP] (wK.south) |- (att.west);
\draw[flowD] (wQ.south) |- (att.east);
\node[xat, minimum width=2.9cm] (cv) at (8.6,-11.65) {$\mathbf{C}^{(D\to P)}=\mathbf{A}\mathbf{V}$};
\draw[flowX] (att.south) -- (cv.north);
\draw[flowP, dsh] (wV.south) |- (cv.west);
\node[oplus] (ad) at (10.75,-11.65) {$+$};
\node[xat, minimum width=1.7cm] (ln) at (12.2,-11.65) {LayerNorm};
\node[ten, fill=XatFill] (xx) at (14.45,-11.65) {$\mathbf{X}\in\mathbb{R}^{R_D\times d}$};
\draw[flowX] (cv) -- (ad);
\draw[flowX] (ad) -- (ln);
\draw[flowX] (ln) -- (xx);
% residual from the Persuadee stream
\draw[flowD, dsh] (tD.east) -- (15.55,-7.15) -- (15.55,-11.05) -- (10.75,-11.05) -- (ad.north);
\node[nt, anchor=south east] at (15.45,-11.05) {residual $\tilde{\mathbf{H}}^{(D)}$};

% ---------------- dual aggregation ------------------------------------------------
% X feeds both readouts: straight down into pooling, and along a bus into the GRU inputs
\coordinate (br) at (14.45,-13.5);
\fill[XatMain] (br) circle (1.6pt);
\draw[XatMain, line width=1.0pt] (br) -- (7.05,-13.5) -- (7.05,-14.75) -- (1.4,-14.75);

\node[ttl, text=TrjBord, anchor=north west] at (0.35,-13.95) {Commitment-trajectory GRU};
\node[trj, minimum width=1.2cm, font=\scriptsize] (g1) at (1.4,-15.6) {GRU};
\node[trj, minimum width=1.2cm, font=\scriptsize] (g2) at (3.4,-15.6) {GRU};
\node[font=\scriptsize] (gd) at (4.75,-15.6) {$\cdots$};
\node[trj, minimum width=1.2cm, font=\scriptsize] (g3) at (6.1,-15.6) {GRU};
\draw[flowT] (g1) -- node[above,font=\tiny]{$\mathbf{h}_1$} (g2);
\draw[flowT] (g2) -- (gd); \draw[flowT] (gd) -- (g3);
\foreach \g/\l in {g1/{\mathbf{x}_1},g2/{\mathbf{x}_2},g3/{\mathbf{x}_{R_D}}}{
  \draw[flowX] (\g.north |- 0,-14.75) -- node[right, font=\tiny, text=Slate, inner sep=1.5pt]{$\l$} (\g.north);}
\node[ten, fill=TrjFill] (hT) at (6.1,-16.55) {$\mathbf{h}_T\in\mathbb{R}^{d}$};
\draw[flowT] (g3.south) -- (hT.north);
\node[nt,anchor=east] at (5.35,-16.55) {read at the last valid Persuadee turn};

\node[ttl, text=TrjBord, anchor=north west] at (8.55,-13.95) {Parallel attention pooling};
\node[trj, minimum width=6.4cm] (alph) at (12.1,-15.1) {$\alpha_t=\operatorname{softmax}_t\!\big(\mathbf{w}^{\top}\tanh(\mathbf{W}\mathbf{x}_t)\big)$};
\node[trj, minimum width=4.1cm] (pv) at (12.1,-16.2) {$\mathbf{p}=\sum_{t}\alpha_t\,\mathbf{x}_t\in\mathbb{R}^{d}$};
\draw[flowT] (alph) -- (pv);
\draw[flowX] (xx.south) -- (xx.south |- alph.north);

% ---------------- fusion, head, output --------------------------------------------
\node[fin,minimum width=5.4cm] (cat) at (8,-17.85) {$\mathbf{z}=[\,\mathbf{h}_T\,\Vert\,\mathbf{p}\,]\in\mathbb{R}^{2d}$};
\draw[flowT] (hT.south) -- (hT.south |- cat.north);
\draw[flowT] (pv.south) |- (cat.east);
\node[fin,minimum width=5.4cm] (head) at (8,-18.95) {Binary head: 2-layer MLP};
\node[fin,minimum width=5.4cm] (yh) at (8,-20.05) {$\hat{y}\in\{\text{yes},\text{no}\}$};
\draw[flow] (cat) -- (head); \draw[flow] (head) -- (yh);
\node[nt,anchor=west, text width=4.4cm, align=left] (lossnote) at (11.45,-19.5)
  {Training loss: class-balanced focal ($\gamma{=}2$, $\beta{=}0.999$) plus $0.30\times$ supervised contrastive ($\tau{=}0.1$)};
\draw[dotted, line width=0.8pt, draw=SlateSoft] (head.east) -- (head.east -| lossnote.west) -- (lossnote.north west);

% ---------------- legend ------------------------------------------------------------
\begin{scope}[shift={(0.5,-21.1)}]
  \node[per, minimum width=0.45cm, minimum height=0.3cm] (l1) at (0,0) {};
  \node[anchor=west, font=\scriptsize] at (0.35,0) {Persuader};
  \node[dec, minimum width=0.45cm, minimum height=0.3cm] at (2.4,0) {};
  \node[anchor=west, font=\scriptsize] at (2.75,0) {Persuadee};
  \node[xat, minimum width=0.45cm, minimum height=0.3cm] at (4.9,0) {};
  \node[anchor=west, font=\scriptsize] at (5.25,0) {Cross-attention};
  \node[trj, minimum width=0.45cm, minimum height=0.3cm] at (8.1,0) {};
  \node[anchor=west, font=\scriptsize] at (8.45,0) {Trajectory / pooling};
  \node[fin,minimum width=0.45cm, minimum height=0.3cm] at (11.7,0) {};
  \node[anchor=west, font=\scriptsize] at (12.05,0) {Fusion / output};
\end{scope}

\end{tikzpicture}
}
  \caption{Detailed neural architecture and tensor specifications of the AntiphonyClassifier, depicting exact layer dimensions, 2-layer independent stream Transformers, directional cross-attention ($\mathbf{Q}_{\mathcal{D}} \mathbf{K}_{\mathcal{P}}^\top / \sqrt{d_k}$), parallel commitment-trajectory GRU and learned attention pooling, and composite Class-Balanced Focal and Supervised Contrastive loss heads.}
  \label{fig:antiphony_arch}
\end{figure}

\begin{summarybox}[Directional Cross-Attention and Trajectory Invariant]
In the \textit{AntiphonyClassifier}, asymmetric conversational roles are embedded directly into the neural wiring:
\begin{equation}
  \mathbf{A}^{(\mathcal{D} \to \mathcal{P})} = \text{softmax}\left(\frac{\mathbf{Q}_{\mathcal{D}} \mathbf{K}_{\mathcal{P}}^\top}{\sqrt{d_k}} + \mathbf{M}_{\mathcal{P}}\right), \quad \mathbf{C}^{(\mathcal{D} \to \mathcal{P})} = \mathbf{A}^{(\mathcal{D} \to \mathcal{P})} \mathbf{V}_{\mathcal{P}}
\end{equation}
The decider query $\mathbf{Q}_{\mathcal{D}}$ selectively attends to the persuader's rhetorical tactics $\mathbf{K}_{\mathcal{P}}, \mathbf{V}_{\mathcal{P}}$, updating the commitment-trajectory state $\mathbf{h}_t = \text{GRU}(\mathbf{x}_t, \mathbf{h}_{t-1})$ read strictly at the dialogue's true final decider turn. This architectural design ensures the model tracks temporal decision momentum while completely ignoring trailing post-decision pleasantries.
\end{summarybox}

\clearpage
\section{Architectural Evolution: The Path to Dual-Stream Antiphony}
\label{sec:evolution}

The AntiphonyClassifier was not designed in isolation; it emerged from a systematic multi-stage empirical investigation across four distinct evolutionary phases, documented in \Cref{fig:architectural_evolution}.

\begin{figure}[htbp]
  \centering
  \resizebox{\textwidth}{!}{% Four architectural stages. Metrics are the best test macro-F1 of each stage
% (153-row test split).
%
% Layout grid (all lengths in cm, y grows downward from the card top at +0.10):
%   card width 3.85, gutter 0.42, inner margin 0.225 -> content width 3.40
%   x origins   : 0 / 4.27 / 8.54 / 12.81      card centre = x + 1.925
%   half width  : 1.58 boxes at x+1.015 and x+2.835, 0.24 apart
%   header band : +0.10 .. -0.86
%   input frame : -1.72 .. -2.16, token chips on its centreline at -1.94;
%                 the encoder arrow leaves the frame edge, not a chip
%   encoder row : -2.78, shared by all four stages so the stacks line up
%   clear gap   : tuned per stage so each stack fills its panel evenly --
%                 0.72 (stage 1, 3 rows), 0.36 (stages 2-3), 0.26 (stage 4,
%                 6 rows); fan-out rows add ~0.14 for the splitter bar
%   output row  : -6.86, shared by the stage 2 and stage 3 heads
%   panel rule  : -8.08    change list : -8.22
%   metric rule : -10.42   bar track   : -10.98 .. -11.26   card foot : -11.50
\begin{tikzpicture}[x=1cm,y=1cm,
  % slimmer heads than the shared style: these connectors are short, and the
  % default 4.2pt tip swallows the whole shaft.
  flow/.append style ={-{Stealth[length=3.4pt,width=3.0pt]}},
  flowP/.append style={-{Stealth[length=3.4pt,width=3.0pt]}},
  flowD/.append style={-{Stealth[length=3.4pt,width=3.0pt]}},
  flowX/.append style={-{Stealth[length=3.4pt,width=3.0pt]}},
  flowT/.append style={-{Stealth[length=3.4pt,width=3.0pt]}},
  % The turn sequence is one object. The encoder arrow leaves this frame, so it
  % reads as "the whole sequence feeds the encoder" rather than appearing to
  % come out of whichever chip happens to sit on the card axis.
  seqbox/.style={rectangle, rounded corners=3pt, draw=NeuBord!45,
                 fill=NeuFill!40, line width=0.6pt}]

% hanging-indent bullet used by the four "what changed" lists
\def\bi{\par\hangindent=0.26cm\hangafter=1\noindent\makebox[0.26cm][l]{{\tiny$\bullet$}}}

% ------------- card frames and headers ---------------------------------------
\foreach \px/\ph/\pn/\pt in {%
    0/64748B/1/{Monolithic Flat},
    4.27/475569/2/{Hierarchical Interleaved},
    8.54/334155/3/{Gated \& Decoupled},
    12.81/0F172A/4/{Dual-Stream Antiphony}}{
  \definecolor{stg}{HTML}{\ph}
  \fill[white, rounded corners=5pt] (\px,0.10) rectangle ++(3.85,-11.60);
  \fill[stg, rounded corners=5pt]   (\px,0.10) rectangle ++(3.85,-0.96);
  \fill[stg]                        (\px,-0.40) rectangle ++(3.85,-0.46);
  \draw[rounded corners=5pt, draw=stg, line width=0.9pt]
        (\px,0.10) rectangle ++(3.85,-11.60);
  \node[text=white!85, font=\tiny\bfseries]       at (\px+1.925,-0.29) {STAGE \pn};
  \node[text=white,    font=\scriptsize\bfseries] at (\px+1.925,-0.62) {\pt};
  % rules that split each card into architecture / changes / metric bands
  \draw[draw=NeuBord!35, line width=0.5pt] (\px+0.225,-8.08)  -- (\px+3.625,-8.08);
  \draw[draw=NeuBord!35, line width=0.5pt] (\px+0.225,-10.42) -- (\px+3.625,-10.42);
}
% stage-to-stage progression arrows, kept inside the gutters
\foreach \px in {3.85,8.12,12.39}{
  \draw[-{Stealth[length=5pt,width=4.2pt]}, line width=1.4pt, draw=Slate]
        (\px+0.05,-0.38) -- (\px+0.37,-0.38);
}

% ------------- descriptors ------------------------------------------------------
\node[anchor=north, text width=3.40cm, align=center, font=\scriptsize, text=SlateSoft]
  at (1.925,-0.98)  {Flat sequence baseline};
\node[anchor=north, text width=3.40cm, align=center, font=\scriptsize, text=SlateSoft]
  at (6.195,-0.98)  {Hierarchical dialogue transformer};
\node[anchor=north, text width=3.40cm, align=center, font=\scriptsize, text=SlateSoft]
  at (10.465,-0.98) {Last-turn gating and specialist routing};
\node[anchor=north, text width=3.40cm, align=center, font=\scriptsize, text=SlateSoft]
  at (14.735,-0.98) {Dual-stream AntiphonyClassifier};

% =============== Stage 1 : flat baseline ========================================
\begin{scope}[shift={(0,0)}]
  % Each caption is centred over the group it names, and the two groups are
  % held 0.24 apart so the truncation rule gets a lane of its own.
  \node[font=\tiny, text=SlateSoft] at (1.295,-1.54) {256 tokens kept};
  \node[font=\tiny, text=PerBord]   at (2.915,-1.54) {ending discarded};
  % only the kept window is framed: that is what actually reaches the encoder
  \draw[seqbox] (0.375,-1.72) rectangle (2.215,-2.16);
  \foreach \i in {0,1,2,3,4}{
    \node[ten, minimum size=0.30cm] at (0.575+0.36*\i,-1.94) {}; }
  \foreach \i in {0,1,2}{
    \node[ten, minimum size=0.30cm, draw=NeuBord!40, dashed]
          at (2.555+0.36*\i,-1.94) {}; }
  \draw[PerMain, dashed, dash pattern=on 2pt off 1.6pt, line width=0.8pt]
        (2.285,-1.66) -- (2.285,-2.22);
  \node[neu, minimum width=3.40cm] (e1)  at (1.925,-2.78) {\scriptsize Encoder, flat input};
  \node[neu, minimum width=3.40cm] (p1)  at (1.925,-4.12) {\scriptsize {[}CLS{]} / mean pooling};
  \node[neu, minimum width=1.58cm] (h1a) at (1.015,-5.54) {\scriptsize Binary};
  \node[neu, minimum width=1.58cm] (h1b) at (2.835,-5.54) {\scriptsize Modifier};
  % leaves the centre of the KEPT frame (1.295), not the card axis: only that
  % window reaches the encoder. The encoder box is full width, so it still lands.
  \draw[flow]   (1.295,-2.16) -- (1.295,-2.47);
  \draw[flow]   (e1) -- (p1);
  \draw[flow,-] (p1.south)    -- (1.925,-4.83);
  \draw[flow,-] (1.015,-4.83) -- (2.835,-4.83);
  \draw[flow]   (1.015,-4.83) -- (h1a.north);
  \draw[flow]   (2.835,-4.83) -- (h1b.north);
\end{scope}

% =============== Stage 2 : hierarchical =========================================
\begin{scope}[shift={(4.27,0)}]
  \draw[seqbox] (0.32,-1.72) rectangle (3.53,-2.16);
  \foreach \i [evaluate=\i as \j using {int(\i+1)}] in {0,1,2,3}{
    \node[ten, minimum width=0.55cm, minimum height=0.30cm]
          at (0.67+0.68*\i,-1.94) {\tiny $u_{\j}$}; }
  \node[font=\tiny, text=SlateSoft] at (3.33,-1.94) {$\cdots$};
  \node[neu, minimum width=3.40cm] (e2)  at (1.925,-2.78) {\scriptsize Utterance encoder, per turn};
  \node[neu, minimum width=3.40cm] (p2)  at (1.925,-3.76) {\scriptsize Turn vectors + role + position};
  \node[amb, minimum width=3.40cm] (d2)  at (1.925,-4.74) {\scriptsize Dialogue transformer, 4 layers};
  \node[amb, minimum width=3.40cm] (a2)  at (1.925,-5.72) {\scriptsize Learned attention pooling};
  \node[neu, minimum width=1.58cm] (h2a) at (1.015,-6.86) {\scriptsize Binary};
  \node[neu, minimum width=1.58cm] (h2b) at (2.835,-6.86) {\scriptsize Modifier};
  \draw[flow]   (1.925,-2.16) -- (e2.north);
  \draw[flow]   (e2) -- (p2);
  \draw[flow]   (p2) -- (d2);
  \draw[flow]   (d2) -- (a2);
  \draw[flow,-] (a2.south)    -- (1.925,-6.35);
  \draw[flow,-] (1.015,-6.35) -- (2.835,-6.35);
  \draw[flow]   (1.015,-6.35) -- (h2a.north);
  \draw[flow]   (2.835,-6.35) -- (h2b.north);
\end{scope}

% =============== Stage 3 : gating and routing ===================================
\begin{scope}[shift={(8.54,0)}]
  \draw[seqbox] (0.32,-1.72) rectangle (3.53,-2.16);
  \foreach \i [evaluate=\i as \j using {int(\i+1)}] in {0,1,2,3}{
    \node[ten, minimum width=0.55cm, minimum height=0.30cm]
          at (0.67+0.68*\i,-1.94) {\tiny $u_{\j}$}; }
  \node[font=\tiny, text=SlateSoft] at (3.33,-1.94) {$\cdots$};
  \node[neu, minimum width=3.40cm] (e3) at (1.925,-2.78) {\scriptsize Utterance encoder, per turn};
  \node[amb, minimum width=3.40cm] (d3) at (1.925,-3.76) {\scriptsize Dialogue transformer, 4 layers};
  \node[amb, minimum width=3.40cm] (a3) at (1.925,-4.74) {\scriptsize Attention pooling};
  \node[oplus] (g3) at (1.925,-5.66) {\scriptsize $\otimes$};
  \node[amb, minimum width=1.25cm, font=\tiny, inner xsep=2pt] (l3) at (3.00,-5.66) {last\\Persuadee};
  \node[font=\tiny, text=SlateSoft, anchor=east] at (1.66,-5.66) {gate};
  % The 0.834 model is the GATED TWO-HEADED one, so both heads are drawn. The
  % modifier head is dashed because decoupling (deleting it) is what closes this
  % stage -- drawing a binary-only head here would attribute 0.834 to a model
  % that never produced it.
  \node[neu, minimum width=1.58cm] (h3a) at (1.015,-6.86) {\scriptsize Binary};
  \node[neu, minimum width=1.58cm, dashed, draw=NeuBord!45, fill=white,
        text=NeuBord!75] (h3b) at (2.835,-6.86) {\scriptsize Modifier};
  \node[font=\tiny, text=PerBord] at (2.835,-7.36) {then dropped};
  \draw[flow] (1.925,-2.16) -- (e3.north);
  \draw[flow] (e3) -- (d3);
  \draw[flow] (d3) -- (a3);
  \draw[flow] (a3.south) -- (g3.north);
  \draw[flow] (l3.west)  -- (g3.east);
  \draw[flow,-] (g3.south)    -- (1.925,-6.35);
  \draw[flow,-] (1.015,-6.35) -- (2.835,-6.35);
  \draw[flow]   (1.015,-6.35) -- (h3a.north);
  \draw[flow, draw=NeuBord!55, dashed, dash pattern=on 2pt off 1.6pt]
        (2.835,-6.35) -- (h3b.north);
\end{scope}

% =============== Stage 4 : Antiphony ============================================
\begin{scope}[shift={(12.81,0)}]
  \draw[seqbox] (0.32,-1.72) rectangle (3.53,-2.16);
  \foreach \i [evaluate=\i as \j using {int(\i+1)}] in {0,2}{
    \node[per, minimum width=0.55cm, minimum height=0.30cm, inner sep=0pt, font=\tiny]
          at (0.67+0.68*\i,-1.94) {$u_{\j}$}; }
  \foreach \i [evaluate=\i as \j using {int(\i+1)}] in {1,3}{
    \node[dec, minimum width=0.55cm, minimum height=0.30cm, inner sep=0pt, font=\tiny]
          at (0.67+0.68*\i,-1.94) {$u_{\j}$}; }
  \node[font=\tiny, text=SlateSoft] at (3.33,-1.94) {$\cdots$};
  \node[neu, minimum width=3.40cm] (e4) at (1.925,-2.78) {\scriptsize Utterance encoder, per turn};
  \node[per, minimum width=1.58cm] (sp) at (1.015,-3.88) {\scriptsize Persuader\\[-2pt]\scriptsize stream};
  \node[dec, minimum width=1.58cm] (sd) at (2.835,-3.88) {\scriptsize Persuadee\\[-2pt]\scriptsize stream};
  \node[xat, minimum width=3.40cm] (cx) at (1.925,-4.92)
       {\scriptsize Directional cross-attention\\[-1.5pt]%
        {\tiny Q\,$\leftarrow$\,Persuadee, K/V\,$\leftarrow$\,Persuader}};
  \node[trj, minimum width=1.58cm] (gr) at (1.015,-5.98) {\scriptsize GRU};
  \node[trj, minimum width=1.58cm] (pl) at (2.835,-5.98) {\scriptsize Pooling};
  \node[oplus, minimum size=0.42cm] (cc) at (1.925,-6.76) {\tiny $[\,;]$};
  \node[fin, minimum width=3.40cm] (h4) at (1.925,-7.54) {\scriptsize Binary head};
  \draw[flow]    (1.925,-2.16) -- (e4.north);
  \draw[flow,-]  (e4.south)    -- (1.925,-3.28);
  \draw[flow,-]  (1.015,-3.28) -- (2.835,-3.28);
  \draw[flowP]   (1.015,-3.28) -- (sp.north);
  \draw[flowD]   (2.835,-3.28) -- (sd.north);
  \draw[flowP]   (sp.south) -- (sp.south |- cx.north);
  \draw[flowD]   (sd.south) -- (sd.south |- cx.north);
  \draw[flowX,-] (cx.south)    -- (1.925,-5.48);
  \draw[flowX,-] (1.015,-5.48) -- (2.835,-5.48);
  \draw[flowX]   (1.015,-5.48) -- (gr.north);
  \draw[flowX]   (2.835,-5.48) -- (pl.north);
  \draw[flowT]   (gr.south) -- (cc.north west);
  \draw[flowT]   (pl.south) -- (cc.north east);
  \draw[flowT]   (cc.south) -- (h4.north);
\end{scope}

% ------------- what changed -------------------------------------------------------
\node[anchor=north west, text width=3.40cm, font=\scriptsize, align=left,
      text=SlateSoft, inner sep=0pt] at (0.225,-8.22)
  {\parskip=1.6pt\parindent=0pt
   \bi 256-token right truncation
   \bi role embeddings, focal loss, marker-safe augmentation
   \bi dual-task joint heads\par};
\node[anchor=north west, text width=3.40cm, font=\scriptsize, align=left,
      text=SlateSoft, inner sep=0pt] at (4.495,-8.22)
  {\parskip=1.6pt\parindent=0pt
   \bi 32 turns $\times$ 64 tokens = 2048-token capacity
   \bi back-translation data augmentation
   \bi tri-backbone soft ensemble\par};
\node[anchor=north west, text width=3.40cm, font=\scriptsize, align=left,
      text=SlateSoft, inner sep=0pt] at (8.765,-8.22)
  {\parskip=1.6pt\parindent=0pt
   \bi last-turn commitment gate
   \bi multi-pivot translation diversity
   \bi confidence specialist routing
   \bi modifier head decoupled\par};
\node[anchor=north west, text width=3.40cm, font=\scriptsize, align=left,
      text=SlateSoft, inner sep=0pt] at (13.035,-8.22)
  {\parskip=1.6pt\parindent=0pt
   \bi role-pure streams
   \bi directional cross-attention
   \bi trajectory GRU plus pooling\par};

% ------------- performance bars ----------------------------------------------------
% track = 3.40 cm spans macro-F1 0.40 to 0.90, so width = (v - 0.40) / 0.50 * 3.40
\foreach \px/\pw/\pv/\pm/\ph in {%
    0/2.040/0.70/{Role-Flat, RoBERTa}/64748B,
    4.27/2.863/0.821/{Aug-Hier, DeBERTa}/475569,
    8.54/2.951/0.834/{Gated-Hier, RoBERTa}/334155,
    12.81/3.142/0.862/{Antiphony, DeBERTa}/0F172A}{
  \definecolor{stg}{HTML}{\ph}
  \node[font=\footnotesize\bfseries, text=Slate, anchor=west] at (\px+0.225,-10.70) {\pv};
  \node[font=\scriptsize, text=SlateSoft, anchor=east]        at (\px+3.625,-10.70) {\pm};
  \fill[NeuFill, rounded corners=2pt] (\px+0.225,-10.98) rectangle ++(3.40,-0.28);
  \fill[stg,     rounded corners=2pt] (\px+0.225,-10.98) rectangle ++(\pw,-0.28);
}

\end{tikzpicture}
}
  \caption{Four-stage architectural evolution of intention classification models, depicting key innovations, critical bug diagnoses, and peak test Macro-F1 across stages.}
  \label{fig:architectural_evolution}
\end{figure}

\subsection{Stage 1: Monolithic Flat Baselines and Forensic Debugging}
\begin{itemize}
  \item \textbf{Flat Sequence Base Model (M1):} Concatenated dialogues truncated to 256 tokens with \texttt{[CLS]} pooling and independent binary and modifier heads. Reached $0.68$ Macro-F1 with RoBERTa, but DeBERTa collapsed to the majority class and modifier heads completely failed on rare classes.
  \item \textbf{Speaker-Aware Flat Encoding (M2):} Added learnable role embeddings ($\text{Persuader}=0, \text{Persuadee}=1$) prior to masked mean-pooling. Slightly improved recall ($0.70$ Macro-F1), but remained bottlenecked by 256-token truncation.
  \item \textbf{Dynamic Augmentation Exploration (M3):} Added EDA text augmentation \cite{wei2019eda} and raw inverse-frequency class weighting. Training completely collapsed. Forensic auditing identified two critical bugs: (i) checkpoint selection was mistakenly scored on \texttt{pos\_label="yes"} F1, allowing an all-yes trivial model to score $1.0$ and be saved, and (ii) raw inverse-frequency weights on the 13-example \texttt{conditional} class caused catastrophic gradient spikes.
  \item \textbf{Metric-Corrected Focal Loss Baseline (M4):} Fixed the evaluation harness to use true unweighted Macro-F1; implemented Class-Balanced Focal Loss ($\beta=0.999$); cascaded the binary prediction into the modifier head; added in-batch Supervised Contrastive Loss. Training stabilized, but performance plateaued at $0.649$ Macro-F1.
  \item \textbf{Token-Safe Augmentation Audit (M5):} An audit revealed that $57.5\%$ of speaker markers (\texttt{[Persuader]}) in raw transcripts were glued to newlines without spaces. Standard whitespace augmentation tokenizers fragmented these tags, corrupting role markers in $15.9\%$ of augmented instances. Implementing regex-based tokenization and Persuadee-biased mutation weights resolved the corruption, but Macro-F1 remained capped ($0.670$), proving that flat 256-token architectures had reached their theoretical ceiling.
\end{itemize}

\subsection{Stage 2: Hierarchical Discourse Modeling}
\begin{itemize}
  \item \textbf{Hierarchical Dialogue Transformer (M6):} The foundational breakthrough: turns encoded individually ($32 \times 64 = 2{,}048$ token capacity) via a shared encoder, pooled into turn vectors, combined with role and positional embeddings, and processed by a 4-layer inter-turn dialogue transformer with learned attention pooling. Macro-F1 surged to \textbf{$0.793$} (RoBERTa), \textbf{$0.797$} (DeBERTa), and \textbf{$0.807$} (TOD-BERT).
  \item \textbf{Flat Back-Translation Control (M7):} Applied French back-translation to the flat M5 model as an experimental control; Macro-F1 dropped to $0.590$--$0.650$, proving that data augmentation cannot overcome architectural truncation.
  \item \textbf{Hierarchical Back-Translation (M8):} Combined the M6 hierarchical architecture with back-translation. Binary Macro-F1 advanced to \textbf{$0.821$} (DeBERTa), but modifier F1 fell, revealing that translation round-trips wash away subtle hedge words.
  \item \textbf{Tri-Model Soft Ensemble (M9):} Soft probability averaging across saved M6 checkpoints (RoBERTa, DeBERTa, TOD-BERT) achieved $0.818$ Macro-F1. 20 stochastic Monte Carlo dropout forward passes demonstrated that prediction variance was significantly lower on correct instances ($0.0220$) than incorrect instances ($0.0297$), providing a calibrated uncertainty signal.
\end{itemize}

\subsection{Stage 3: Recency Gating and Decoupled Baselines}
\begin{itemize}
  \item \textbf{Last-Persuadee-Turn Gating (M10):} Added a sigmoid gate explicitly anchoring the pooled representation to the final Persuadee hidden state, initialized with a $+3.0$ bias to prevent early deflection. Achieved \textbf{$0.834$} Macro-F1 with RoBERTa, establishing the peak single-model binary score of the two-headed lineage.
  \item \textbf{Multi-Pivot Back-Translation Diversity (M11):} Implemented French and German back-translation under beam search and random sampling; DeBERTa reached \textbf{$0.829$}, but modifier F1 continued to degrade.
  \item \textbf{Specialist Confidence Router (M12):} Decoupled the tasks without retraining: routing binary classification to M10 Gated RoBERTa ($0.834$) and modifier classification to M6 Plain DeBERTa ($0.452$).
  \item \textbf{Decoupled Single-Stream Baseline:} Deleted the modifier head entirely, dedicating all capacity to a 2-layer binary head. Reached \textbf{$0.825$} on second-pass labels and \textbf{$0.8359$} on third-pass labels.
\end{itemize}

\subsection{Stage 4: Role-Pure Dual-Stream Antiphony Architecture}
The Dual-Stream Antiphony architecture implemented speaker-disentangled cross-attention, establishing project-wide peak performance ($0.8623$ on the second-pass benchmark, $0.8446$ on the strict third-pass benchmark), definitively proving the necessity of modeling asymmetric conversational dynamics.

\section{Design Refinement and Simulation}
\label{sec:detailed-design}

\subsection{Capacity Calculations and Sequence Length Budgeting}
In conversational modeling, sequence capacity must balance linguistic coverage against quadratic attention memory scaling ($\mathcal{O}(N^2)$). In the P4G corpus, dialogue length exhibits a median of 344 words across an average of 20 turns, with the 95th percentile reaching 23 turns. Flat transformers capped at 256 tokens cover fewer than 180 words, truncating over $50\%$ of conversational content. 

In Antiphony, the sequence budget allocates $U=32$ turns with $L=64$ tokens per turn, providing a total token capacity of:
\begin{equation}
  C_{\text{tokens}} = 32 \times 64 = 2{,}048 \text{ tokens}
\end{equation}
This budget covers $99.7\%$ of all conversational tokens across the corpus without truncation. For role-pure streams, setting $R=20$ turns per role provides capacity for up to 40 total turns, covering $98.8\%$ of dialogue trajectories.

\subsection{Memory Profiling and Hardware Guards}
Processing a batch of size $B$ with $U$ turns of length $L$ through a 768-dimensional transformer requires flattening to $[B \cdot U, L]$. For $B=4$ and $U=32$, this requires encoding 128 sequences of length 64 in a single forward pass. On an NVIDIA Tesla T4 GPU (16\,GB VRAM), peak activation memory for DeBERTa-v3 under standard precision exceeds hardware limits. We implemented three strict hardware guards:
\begin{enumerate}
  \item \textbf{Micro-Batch Slicing:} Training batch size is set to $B=2$, with evaluation batch size $B=4$.
  \item \textbf{Gradient Checkpointing:} Intermediate activation maps are discarded during the forward pass and recomputed during backpropagation, reducing activation memory by $\approx 60\%$.
  \item \textbf{PyTorch Dynamic Memory Allocation:} Enabled \texttt{expandable\_segments:True} in the PyTorch CUDA allocator to eliminate memory fragmentation during variable-length sequence packing.
\end{enumerate}

\section{Application of Tools}
\label{sec:tools}

The project utilized a cohesive suite of specialized computational tools:
\begin{itemize}
  \item \textbf{PyTorch 2.x \cite{paszke2019pytorch}:} Provided core automatic differentiation, custom CUDA autograd layers, and modular neural network implementations.
  \item \textbf{Hugging Face Transformers \cite{wolf2020transformers}:} Supplied pre-trained transformer backbones (RoBERTa, DeBERTa-v3, TOD-BERT, ELECTRA) and tokenization engines.
  \item \textbf{XGBoost \cite{chen2016xgboost} and Scikit-Learn:} Utilized for classical baseline construction, tabular stacking experiments, calibration curve fitting, and bootstrap evaluation metrics.
  \item \textbf{OPUS-MT \cite{tiedemann2020opusmt}:} Integrated for multilingual back-translation pipelines across French and German pivots.
  \item \textbf{Anthropic Claude 5 Sonnet:} Utilized as the automated frontier LLM rater for contextual multi-label strategy annotation under human calibration checkpoints.
  \item \textbf{Google Gemini API (\texttt{gemini-3.5-flash-lite}):} Employed for organic paraphrasing, controlled jitter injection, and synthetic dialogue auditing.
  \item \textbf{MiKTeX and PGF/TikZ:} Leveraged for publication-grade mathematical typesetting and native vector architectural diagramming.
\end{itemize}
\textbf{Tool Limitations:} Pre-trained transformer checkpoints inherit tokenization biases and vocabulary constraints; free-tier cloud environments enforce 12-hour timeout limits; and proprietary LLM APIs introduce rate-limit constraints and network latency.

\section{Implementation}
\label{sec:implementation}

The development process adhered to modular software engineering practices. All code was organized into self-contained, reproducible Jupyter notebooks paired with comprehensive markdown architectural summaries (\texttt{architecture-v*.md}). 

\subsection{Version Control and Development History}
The codebase was tracked using Git \cite{git}. The repository logs 14 verified commits spanning September 5 to September 20, 2026, with primary engineering contributions from Abhishek Roy (9 commits), Nawriz Ahmed Turjo (3 commits), and Shams Hossain Simanto (2 commits). Earlier exploratory development was executed in dated local checkpoint notebooks.

\subsection{Deterministic Seeding and Audit Artifacts}
To guarantee reproducible outcomes, all scripts execute a centralized seeding utility setting Python \texttt{random}, NumPy, and PyTorch CPU/CUDA RNG states to Seed $= 42$. All trained model weights, out-of-fold validation predictions, and test evaluation matrices were serialized as immutable \texttt{.pt} and \texttt{.jsonl} artifacts.

\section{Deployment}
\label{sec:deployment}

\subsection{Prototype Execution Pipeline}
In its present form, the system operates as an offline research prototype executed within Kaggle and Python environments. The end-to-end inference execution flow for an unlabelled dialogue proceeds as follows:
\begin{equation}
  \text{Transcript } \mathcal{D} \xrightarrow{\text{Tokenize}} [U \times 64] \xrightarrow{\text{Encoder}} \mathbf{U} \xrightarrow{\text{Split \& Cross-Attend}} \mathbf{X} \xrightarrow{\text{GRU + Pool}} \mathbf{z} \xrightarrow{\text{MLP}} \hat{y}
\end{equation}

\subsection{Resource Consumption}
Inference on a complete 20-turn dialogue executes in approximately $45$\,ms on an NVIDIA T4 GPU ($380$\,ms on an 8-core Intel CPU under batch size 1). Measured GPU training runtimes across model variants: Metric-Corrected Baseline ($\approx 20$\,min), Hierarchical Transformer ($\approx 2.5$\,h), Hierarchical Back-Translation ($\approx 5.0$\,h), Last-Turn Gated Model ($\approx 2.5$\,h), Decoupled Single-Stream ($\approx 3.5$\,h), and Antiphony ($\approx 3.0$\,h). The category clean rerun took 177 minutes, and strategy ensemble search required $\approx 4.0$\,hours.

\subsection{Operational Boundaries and Omissions}
We state plainly that this project is an academic research prototype. \textbf{No commercial production serving microservice, optimized runtime graph export, model quantization, automated deployment pipeline, or live latency SLA benchmarks were constructed.} Deploying the system into production would require implementing streaming turn ingestion, model quantization, and robust fail-soft API infrastructure.